% ======================================================================
% Chapter 17: Quantum Chromodynamics
% (book pages 545-597; PDF pages 568-621)
% ======================================================================
\chapter{Quantum Chromodynamics}

The key to constructing a realistic model of the strong interactions is the
phenomenon of asymptotic freedom.  Chapter~14 described the experimental
discovery of this phenomenon, while Chapter~16 presented the theoretical
proof that non-Abelian gauge theories are asymptotically free.  We are now
ready to explore the consequences of these discoveries.

We will begin by arguing that the most natural candidate for a model of the
strong interactions is the non-Abelian gauge theory with gauge group $SU(3)$,
coupled to fermions (quarks) in the fundamental representation.  This theory
is known as Quantum Chromodynamics, or QCD.  After some general discussion
of QCD in Section~17.1, we will investigate a number of specific QCD
scattering processes in Sections~17.2 through~17.4.  The most interesting
application of QCD, however, is of a somewhat more sophisticated nature; it
comes in the prediction of a pattern of slow violations of the Bjorken scaling
relation discussed in Chapter~14.  Section~17.5 develops the additional
theoretical tools that are needed to understand these violations.

Although this chapter includes many references to experiments, we remind the
reader that, for QCD as for QED or critical phenomena, this book is primarily
a textbook of theoretical methods rather than a review and interpretation of
experimental data.  The details of experimental techniques and results on
strong interaction physics are reviewed in a number of excellent texts (see
the bibliography).  We hope that this chapter will give the theoretical
foundation necessary to illuminate and interpret these results.

\section{From Quarks to QCD}

Our current theoretical picture of the strong interactions began with the
identification of the elementary fermions that make up the proton and other
hadrons.  As the properties of these fermions became better understood, the
nature of their interactions became tightly constrained, in a way that led
eventually to a unique candidate theory.  In order to appreciate the
uniqueness of this theory, we begin this chapter with a simplified history of
how it arose.

In 1963, Gell-Mann and Zweig proposed a model that explained the spectrum of
strongly interacting particles in terms of elementary constituents called
\emph{quarks}.  Mesons were expected to be quark-antiquark bound states.
Indeed, the lightest mesons have just the correct quantum numbers to justify
this interpretation; they are spin-0 and spin-1 states of odd parity, just as
we found for fermion-antifermion bound states of zero orbital angular
momentum in Chapter~3.  Baryons were interpreted as bound states of three
quarks.  To explain the electric charges and other quantum numbers of
hadrons, Gell-Mann and Zweig needed to assume three species of quarks, up
($u$), down ($d$), and strange ($s$).  Additional hadrons discovered since
that time require the existence of three more species: charm ($c$), bottom
($b$), and top ($t$).  To make baryons with integer charges, the quarks
needed to be assigned fractional electric charge: $+2/3$ for $u$, $c$, $t$,
and $-1/3$ for $d$, $s$, $b$.  Then, for example, the proton would be a bound
state of $uud$, while the neutron would be a bound state of $udd$.  The six
types of quarks are conventionally referred to as \emph{flavors}.

The quark model had great success in predicting new hadronic states, and in
explaining the strengths of electromagnetic and weak-interaction transitions
among different hadrons.  In particular, the quark model naturally
incorporates the most important symmetry relations among strongly interacting
particles.  If one assumes that the $u$ and $d$ quarks have identical masses
and interactions, the $SU(2)$ group that acts as a unitary rotation of $u$ and
$d$ states,
\begin{equation}
  \begin{pmatrix}u\\ d\end{pmatrix} \to U\begin{pmatrix}u\\ d\end{pmatrix},
  \label{eq:17.1}
\end{equation}
should be a symmetry of the strong interactions.  Indeed, both in nuclear and
in elementary particle physics, the quantum states form multiplets of this
$SU(2)$ symmetry, called \emph{isotopic spin} or \emph{isospin}.  Similarly,
since the strange quark is only a little heavier than the $u$ and $d$ quarks,
it makes sense to consider the symmetry of unitary transformations of the
triplet $(u, d, s)$.  Gell-Mann and Ne'eman showed that the elementary
particles naturally fill out irreducible representations of this $SU(3)$
symmetry.

Despite the phenomenological success of the original quark model, it had two
serious problems.  First, despite considerable effort, free particles with
fractional charge could not be found.  Second, the spectrum of baryons
required the assumption that the wavefunction of the three quarks be totally
symmetric under the interchange of the quark spin and flavor quantum numbers,
contradicting the expectation that quarks, which must have spin $1/2$, should
obey Fermi-Dirac statistics.  The need for this symmetry is most clearly
illustrated in the fact that one of the lightest excited states of the nucleon
is a spin-$3/2$ particle with charge $+2$, the $\Delta^{++}$.  This particle is
readily interpreted as a $uuu$ bound state with zero orbital angular momentum
and all three quark spins parallel.

To reconcile the baryon spectrum with the spin-statistics theorem, Han and
Nambu, Greenberg, and Gell-Mann proposed that quarks carry an additional,
unobserved quantum number, called \emph{color}.  They introduced the ad hoc
assumption that baryon wavefunctions must be totally antisymmetric in color
quantum numbers.  Then, if the quark wavefunctions are totally symmetric in
spin and flavor, they are totally antisymmetric overall, in agreement with
Fermi-Dirac statistics.  The simplest model of color would be to assign quarks
to the fundamental representation of a new, internal $SU(3)$ global symmetry.
Suppressing for a moment the spin and flavor quantum numbers, we can
represent quarks by $q^i$, where $i = 1, 2, 3$ is the color index.  Thus
quarks transform under the fundamental, or ``3'', representation of the color
$SU(3)$ symmetry.  Antiquarks, $q_i$, transform in the $\bar{3}$
representation.  The inner product of a 3 and a $\bar{3}$ is an invariant of
$SU(3)$.  One can also form an invariant by using the totally antisymmetric
combination of three 3's, $\epsilon^{ijk}$: This object transforms under a
unitary transformation according to
\begin{equation}
  \epsilon^{i'j'k'}U_{ii'}U_{jj'}U_{kk'} = (\det U)\epsilon^{ijk} ,
  \label{eq:17.2}
\end{equation}
and so is invariant under $SU(3)$ transformations, which have $\det U = 1$.
Under the postulate that all hadron wavefunctions must be invariant under
$SU(3)$ symmetry transformations, these two types of combinations are the
only simple ones allowed:
\begin{equation}
  \bar{q}_i q^i,\qquad \epsilon_{ijk}q^i q^j q^k,\qquad
  \epsilon^{ijk}\bar{q}_i\bar{q}_j\bar{q}_k .
  \label{eq:17.3}
\end{equation}
That is, the assumption that physical hadrons are singlets under color
implies that the only possible light hadrons are the mesons, baryons, and
antibaryons.

Like the original quark model, the color hypothesis was phenomenologically
successful but raised additional questions: Why should quarks have this
seemingly superfluous property, and what mechanism insures that all hadron
wavefunctions are color singlets?  The answers to these questions came not
from hadron spectroscopy, but from the deep-inelastic scattering experiments
described in Chapter~14 and the ensuing search for a theory of parton binding
with the property of asymptotic freedom.  When it was discovered that
non-Abelian gauge theories have this property, all that remained was to
identify the correct gauge group and fermion representation.  Since the color
symmetry had no other obvious physical role, it was natural to identify this
symmetry with the gauge group, with the colors being the gauge quantum
numbers of the quarks.  This reasoning resulted in a model of the strong
interactions as a system of quarks, of the various flavors, each assigned to
the fundamental representation of the local gauge group $SU(3)$.  The quanta
of the $SU(3)$ gauge field are called \emph{gluons}, and the theory is known
as Quantum Chromodynamics, or QCD.

In this book, we will mainly investigate the properties of QCD in the
high-energy regime, where the coupling constant has become small.  However,
we should point out that one can also study QCD in the regime of strong
coupling, using an approximation scheme introduced by Wilson in which the
continuum gauge theory is replaced by a discrete statistical mechanical
system on a four-dimensional Euclidean lattice.  Using this approximation,
Wilson showed that, for sufficiently strong coupling, QCD exhibits
\emph{confinement} of color: The only finite-energy asymptotic states of the
theory are those that are singlets of color $SU(3)$.  Thus the ad hoc
assumption that explains the spectrum of hadrons turns out to be a
consequence of the non-Abelian gauge theory coupling to color.

\begin{figure}[t]
  \centering
  \includegraphics[width=0.5\textwidth]{fig17.1.png}
  \caption{Gauge electric field configuration associated with the separation
    of color sources in a strong-coupling gauge theory.}
  \label{fig:17.1}
\end{figure}

If one attempts to separate a color-singlet state into colored components---for
example, to dissociate a meson into a quark and an antiquark---a tube of
gauge field forms between the two sources, as shown in Fig.~\ref{fig:17.1}.
In a non-Abelian gauge theory with sufficiently strong coupling, this tube has
fixed radius and energy density, so the energy cost of separating color
sources grows proportionally to the separation.  A force law of this type can
consistently be weak at short distances and strong at long distances,
accounting for the fact that isolated quarks are not observed.  We will
discuss the large-distance, strong-coupling limit of gauge theories further in
the Epilogue.

The short-distance limit of Quantum Chromodynamics can be readily studied
using the Feynman diagram technology that we have developed in previous
chapters.  Here asymptotic freedom makes the coupling weak, and there is a
sensible diagrammatic perturbation theory that begins from the model of free
quarks and gluons.  The following sections treat the elementary interactions
among quarks and gluons that can be observed in high-energy experiments.

\section{$e^{+}e^{-}$ Annihilation into Hadrons}

The simplest reaction involving quarks is the production of quark pairs in
$e^{+}e^{-}$ annihilation, a process that we treated already in Section~5.1.
There we analyzed this process only at the most elementary level, viewing it
as a pure QED reaction in which free quarks are created by a virtual photon.
The diagram for this lowest-order process is shown in
Fig.~\ref{fig:17.2}(a).  The computation of the total cross section includes
a sum over the various color states of the quark fields, and so provides a
confirmation that the number of allowed colors is 3.  Combining the color
factor with the square of the quark electric charges, we found
(Eq.~\eqref{eq:5.16})
\begin{equation}
  \sigma(e^{+}e^{-}\to\text{hadrons}) = \sigma_0\cdot 3\cdot
    \sum_f Q_f^2 ,
  \label{eq:17.4}
\end{equation}
where $\sigma_0$ is the QED cross section for $e^{+}e^{-}\to\mu^{+}\mu^{-}$,
\begin{equation}
  \sigma_0 = \frac{4\pi\alpha^2}{3s} ,
  \label{eq:17.5}
\end{equation}
and the sum in \eqref{eq:17.4} is taken over quark flavors.  This formula
assumes that the center of mass energy is high enough that we can ignore the
quark masses.

\begin{figure}[t]
  \centering
  \includegraphics[width=0.5\textwidth]{fig17.2.png}
  \caption{Diagrams contributing to the process $e^{+}e^{-}\to$ hadrons in
    QCD: (a) the leading-order diagram; (b) corrections of order $\alpha_s$.}
  \label{fig:17.2}
\end{figure}

When we couple the quarks to an $SU(3)$ gauge theory, we add many important
processes that affect both the value of this cross section and the final
states that it includes.  Some of the most important effects cannot be
discussed within the context of perturbation theory.  In particular, though
the leading diagram contains free quarks, the particles that emerge from the
reaction are color-singlet mesons and baryons.  However, we will find that
QCD perturbation theory with quarks and gluons does make a number of
important predictions for $e^{+}e^{-}$ annihilation to hadrons.  The ideas
that we develop in working out these predictions will also apply to many
other strong-interaction processes.

\subsection{Total Cross Section}

The leading corrections to the rate of $e^{+}e^{-}$ annihilation due to gluon
exchange and emission are shown in Fig.~\ref{fig:17.2}(b).  These are
precisely the diagrams computed in the Final Project of Part~I.  The first
two diagrams give a cross section of order $g^2$, where $g$ is the $SU(3)$
gauge coupling, to produce a gluon in addition to the quark and antiquark.
The third diagram must be summed in the amplitude with the leading diagram
to produce a correction to the rate of $q\bar{q}$ production without gluon
emission.  In Part~I, we computed these two contributions as if the strong
interactions were an Abelian gauge theory.  To obtain the corresponding
expressions in QCD, we need only multiply the Abelian formulae by the group
theory factor
\begin{equation}
  \mathrm{tr}[t^at^a] = C_2(r)\cdot\mathrm{tr}[\mathbb{1}]
  = \frac{4}{3}\cdot 3 ,
  \label{eq:17.6}
\end{equation}
where we have used Eq.~(15.97) to evaluate $C_2(r)$ for the
fundamental representation of $SU(3)$.  The factor of 3 is the same color sum
that appears in Eq.~\eqref{eq:17.4}.  Thus we can obtain the correct formulae
for QCD from those of the Final Project by making the replacement
\begin{equation}
  g^2 \to \frac{4}{3}g^2,\qquad \text{or}\qquad
  \alpha_g \to \frac{4}{3}\alpha_s ,
  \label{eq:17.7}
\end{equation}
where
\begin{equation}
  \alpha_s = \frac{g^2}{4\pi}
  \label{eq:17.8}
\end{equation}
is the strong-interaction analogue of the fine-structure constant.

The end result of the Final Project of Part~I was a formula for the total
cross section to produce hadrons in $e^{+}e^{-}$ annihilation.  If we replace
$\alpha_g$ with $(4/3)\alpha_s$, that result becomes
\begin{equation}
  \sigma(e^{+}e^{-}\to\text{hadrons}) = \sigma_0\cdot 3\cdot
    \sum_f Q_f^2\ \Bigl[ 1 + \frac{\alpha_s}{\pi} + \mathcal{O}(\alpha_s^2)
      \Bigr] .
  \label{eq:17.9}
\end{equation}
This is actually the sum of the rates for two elementary processes,
$e^{+}e^{-}\to q\bar{q}$ (including the correction from the third diagram of
Fig.~\ref{fig:17.2}(b)) and $e^{+}e^{-}\to q\bar{q}g$.  Although the rate for
each of these processes is divergent as the gluon mass is taken to zero, that
divergence cancels when they are combined.  This is another example of the
phenomenon of infrared divergence cancellations that we studied for the
example of electron scattering in Sections~6.4 and~6.5.  There we showed
that the dressing of the final state by the emission of soft and collinear
photons does not affect the overall scattering rate.  Here, we see again that
infrared divergences cancel in the total rate, although the sum over real and
virtual gluon corrections leaves over a simple numerical correction.

It is not difficult to understand the cancellation of infrared logarithms
intuitively.  The original process $e^{+}e^{-}\to q\bar{q}$ is extremely
rapid: Since the virtual photon is off-shell by an amount $q^2 = s$, the
quarks are created in a time $1/\sqrt{s}$.  However, the emission of collinear
gluons, and the virtual corrections associated with the exchange of soft
gluons, occur over a much longer time scale.  In the diagrams with gluon
emission, the virtual quark or antiquark is off-mass-shell by an amount
$p^2_{\perp g}$, where $p_{\perp g}$ is the transverse momentum of the gluon
relative to the $q\bar{q}$ system.  Thus this virtual state survives for a
time $1/p_{\perp g}$ before it decays.  Such a slow process cannot affect the
probability that a $q\bar{q}$ pair was produced; it can only affect the
properties of the final state into which the $q\bar{q}$ system will evolve.
By this logic, the only perturbative corrections that can affect the total
cross section are those for which $p_{\perp g} \sim \sqrt{s}$.

Another way to express this conclusion would be to argue that, after
contributions from the infrared-sensitive regions have canceled, the
contributions that remain come from the region of large real or virtual gluon
momenta.  By either argument, formula \eqref{eq:17.9} should be a meaningful
prediction of QCD perturbation theory, even though it involves an integral
over the region of soft gluon emission.

\subsection{The Running of $\alpha_s$}

The formula \eqref{eq:17.9} involves the renormalized coupling $\alpha_s$,
which depends on the renormalization scale $M$.  To obtain a prediction that
is independent of this arbitrary scale, we use the renormalization group.
The $\beta$ function of QCD was computed in Chapter~16:
\begin{equation}
  \beta(g) = -g\,\frac{b_0}{16\pi^2}\,g^2 ,
  \label{eq:17.11}
\end{equation}
where
\begin{equation}
  b_0 = 11 - \frac{2}{3}n_f ,
  \label{eq:17.12}
\end{equation}
for $SU(3)$ with $n_f$ flavors of quarks.  The solution of the renormalization
group equation
\begin{equation}
  \frac{d\alpha_s}{d\log(Q/M)} = \beta(\alpha_s)
\end{equation}
is
\begin{equation}
  \alpha_s(Q) = \frac{\alpha_s(M)}{1 + (b_0\alpha_s/2\pi)\log(Q/M)} .
  \label{eq:17.14}
\end{equation}
The explicit dependence of $\sigma$ on $\alpha_s$ can be found by matching the
successive terms in the expansion of $f(\alpha_s(\sqrt{s}))$ to the terms in
the perturbative expansion.  To the order of the first corrections, we find
simply
\begin{equation}
  \sigma = \sigma_0\cdot 3\cdot \sum_f Q_f^2\ \Bigl[ 1 + \frac{\alpha_s
    (\sqrt{s})}{\pi} + \cdots \Bigr] .
  \label{eq:17.15}
\end{equation}
Thus the Callan-Symanzik equation instructs us to replace the fixed
renormalized coupling $\alpha_s$ with the running coupling constant
$\alpha_s(Q)$, evaluated at $Q^2 = s$.

Because the fixed coupling $\alpha_s$ depends on the arbitrary renormalization
point $M$, it is sometimes useful to remove it from our formulae completely.
To do this, we define a mass scale conventionally called $\Lambda$ (not to be
confused with an ultraviolet cutoff!) satisfying
\begin{equation}
  1 = \bar{g}^2(b_0/8\pi^2)\log(M/\Lambda) .
  \label{eq:17.16}
\end{equation}
Then Eq.~\eqref{eq:17.14} can be rearranged into the form
\begin{equation}
  \alpha_s(Q) = \frac{4\pi}{b_0\log(Q^2/\Lambda^2)} .
  \label{eq:17.17}
\end{equation}
This formula is the clearest expression of the statement that $\alpha_s(Q)$
becomes small as $(\log Q)^{-1}$ for large $Q$.  The momentum scale $\Lambda$
is the scale at which $\alpha_s$ becomes strong as $Q^2$ is decreased.

Experimental measurements of the rate of this reaction and others yield a
value of $\Lambda \approx 200$ MeV.  QCD perturbation theory is valid only
when $Q$ is somewhat larger than this, say above $Q = 1$ GeV, where
$\alpha_s(Q) \sim 0.4$.  The strong interactions become strong at distances
larger than $\sim 1/\Lambda$, which is roughly the size of the light hadrons.

Although the example of the $e^{+}e^{-}$ annihilation cross section is
especially simple, since it depends on only one momentum invariant, similar
conclusions carry over to other QCD predictions.  In analyzing
strong-interaction processes that are sensitive to the quark and gluon
substructure, we will find leading-order formulae for the reaction cross
sections that depend on the renormalized coupling $\alpha_s$.  To make these
expressions satisfy the Callan-Symanzik equation, we must replace this fixed
coupling with the running coupling constant $\alpha_s(Q)$, evaluated with $Q$
of the order of the momentum invariants of the reaction.  Since the running
coupling constant depends only logarithmically on $Q$, we need not worry
about choosing $Q$ precisely.  If we guess the proper scale incorrectly by a
factor of 2, this induces an error in $\alpha_s(Q)$ that is of order
$(\log Q)^{-2} \sim \alpha_s^2(Q)$.  Conversely, this ambiguity would be
resolved by computing to the next order in $\alpha_s$.

Before concluding this formal treatment of the $e^{+}e^{-}$ annihilation
cross section, we should add one qualification.  At the beginning of
Section~12.2, we remarked that renormalization group predictions can be
complicated by the appearance of physical thresholds and their associated
singularities, and so we stated these predictions only for when the relevant
momentum invariant $P^2$ was large and spacelike.  In the present chapter, we
will be concerned with cross sections for quark and gluon reactions, evaluated
on-shell.  This introduces additional complications of principle.  For
example, in order to apply the Callan-Symanzik equation to $\sigma(s)$, we
needed to know that this quantity contains no infrared divergences whose
regulator might provide another dimensionful parameter.  Throughout this
chapter we will assume that similar cancellations of divergences associated
with soft and collinear gluons occur in the processes of interest to us.  The
complete proof of these cancellations in QCD can be carried through, but it
is rather technical.%
\footnote{For a review of the theorems justifying the formulae of perturbative
QCD, see J.~C.~Collins and D.~E.~Soper, \emph{Ann.\ Rev.\ Nucl.\ Sci.}
\textbf{37}, 383 (1987).}  In some cases, an alternative point of view is
possible; one can justify the use of the renormalization group to analyze an
on-shell amplitude by relating it to Green's functions evaluated in the
spacelike region.  This method of analysis, which brings its own insights,
will be the main subject of Chapter~18.

\subsection{Gluon Emission and Jet Production}

A second result of the Final Project of Part~I was a formula for the
differential cross section for $q\bar{q}$ production with gluon emission.
Transcribing this formula to QCD using \eqref{eq:17.7} gives the following
result: Let $x_1$, $x_2$, $x_3$ be the ratios of the quark, antiquark, and
gluon energies to the electron beam energy.  These satisfy $0 < x_i < 1$ and
$x_1 + x_2 + x_3 = 2$.  Then the cross section for $e^{+}e^{-}\to q\bar{q}g$
is given by
\begin{equation}
  \frac{d^2\sigma}{dx_1\,dx_2}(e^{+}e^{-}\to q\bar{q}g)
  = \sigma_0\cdot\Bigl( 3\sum_f Q_f^2 \Bigr)\cdot
    \frac{2\alpha_s}{3\pi}\ \frac{x_1^2+x_2^2}{(1-x_1)(1-x_2)} .
  \label{eq:17.18}
\end{equation}
This cross section is singular as $x_1$ or $x_2$ approaches 1.  The limit
$x_1\to 1$ corresponds to configurations in which the quark has the maximum
possible energy, while the antiquark and the gluon go off in the opposite
direction, sharing the remaining energy.  Then the antiquark and gluon have
almost collinear lightlike momentum vectors and so form a system of very
small invariant mass.  Similarly, the limit $x_2\to 1$ corresponds to
configurations in which the quark and gluon are collinear.  These
singularities are responsible for the divergence of the integrated cross
section in the limit of vanishing gluon mass.

How should we interpret these singularities?  In our general treatment of
bremsstrahlung in Section~6.1, we saw that the emission of a photon by a
scattered electron is enhanced, for collinear radiation, by a factor of order
$\log(s/m^2)$, where $m$ is the mass of the electron.  Thus the total rate
for emitting a collinear photon formally diverges in the limit of zero mass.
The same conclusion holds for the emission of gluons by quarks.  A divergence
that appears for collinear emissions in the limit of zero mass is called a
\emph{mass singularity}.  In QED, we saw that the mass singularity signals a
real physical effect of strong collinear radiation when $q^2$ is large.  In
QCD, we might expect strong gluon radiation in this limit, but we must think
carefully about how this radiation appears experimentally.  Whether a
collinear gluon is radiated or not, the quark and antiquark emerging from the
reaction will undergo further soft interactions with the other products.
These processes must continue, producing quark-antiquark pairs and emitting
and absorbing gluons, until all colored particles are collected into
color-singlet hadrons.  Thus the presence of one or more collinear gluons will
have no noticeable effect on the final state, which consists of two
back-to-back jets of hadrons.  For this reason, formula \eqref{eq:17.18} is
of no use when the gluon transverse momentum is less than the typical scale
of soft gluon interactions, roughly 1 GeV.

When the gluon is emitted with substantial transverse momentum with respect
to the $q\bar{q}$ axis, however, it is not possible for subsequent soft
exchanges to recall or reverse this transverse momentum.  In this case, the
$q\bar{q}g$ system evolves into a system of three distinct jets of hadrons.
Thus, sufficiently far from the collinear regions, we can interpret
Eq.~\eqref{eq:17.18} as the cross section for producing events with three
hadronic jets having energies $x_1$, $x_2$, $x_3$ times the electron beam
energy.

By an analysis similar to that given above for the total cross section, we
can improve Eq.~\eqref{eq:17.18} by replacing the fixed coupling constant
$\alpha_s$ with a running $\alpha_s(Q)$.  A reasonable choice for $Q$ is the
transverse momentum of the gluon, $p_{\perp g}$, described below
Eq.~\eqref{eq:17.9}.  If this transverse momentum is too small, however,
$\alpha_s(Q)$ will be large, and our leading-order formula will break down.
This is a second reason why we cannot use formula \eqref{eq:17.18} when the
transverse momentum of the gluon is less than about 1 GeV.

On the other hand, when the gluon transverse momentum is much larger than
1 GeV, there is no reason to distrust QCD perturbation theory.  Soft processes
cannot disturb the three-jet nature of the hadronic state, and asymptotic
freedom insures that the coupling constant is small, so that the leading
order of perturbation theory will be a good approximation.

The three-jet cross section \eqref{eq:17.18} is a good example of the type of
prediction that one obtains from the use of perturbation theory in QCD.  We
describe a strong-interaction process by the invariant momentum transfer $Q$
it gives to hadronic constituents.  QCD perturbation theory makes predictions
about the flow of energy and momentum in such a reaction into the final
system of jets of hadrons.  If $Q$ is small, perturbation theory is invalid,
and we obtain no useful prediction.  However, if $Q$ is large, the asymptotic
freedom of QCD implies that Feynman diagrams for quarks and gluons will
correctly predict the behavior of the final system of hadronic jets.

\section{Deep Inelastic Scattering}

After $e^{+}e^{-}$ annihilation into hadrons, the next simplest reaction
involving strongly interacting particles is electron scattering from a proton,
or from some other hadron.  At the most elementary level, this reaction can
be viewed as the electromagnetic scattering of an electron from a quark inside
the proton.%
\footnote{The electron could just as well be a muon; all the same formulae
apply in this case.  Leptons can also scatter from quarks via the
neutral-current weak interaction, as we will see in Chapter~20.}  A way to
visualize the process is shown in Fig.~\ref{fig:17.3}.  Call the proton
momentum $P$, and the initial quark momentum $p$.  Call the initial and final
momenta of the electron $k$ and $k'$.  If the final electron momentum is
measured, one can deduce the momentum $q = k-k'$ transferred by the virtual
photon to the hadronic system.  The vector $q$ is spacelike, and one
conventionally writes
\begin{equation}
  q^2 = -Q^2 .
\end{equation}
If the invariant momentum transfer $Q^2$ is large, the quark is ejected from
the proton in a manner that cannot be balanced by subsequent soft processes.
These soft processes will, however, create gluons and quark-antiquark pairs
that eventually neutralize the color and cause the struck quark to materialize
as a jet of hadrons in the direction of the momentum transfer from the
electron.  Typically the total invariant mass of the final hadronic system is
large, since the struck quark carries a large momentum with respect to the
other ``spectator'' quarks.  In this case, the process is referred to as
\emph{deep inelastic scattering}.

\begin{figure}[t]
  \centering
  \includegraphics[width=0.5\textwidth]{fig17.3.png}
  \caption{Deep inelastic scattering in QCD.  The diagram shows the flow of
    momentum when a high-energy electron scatters from a quark taken from the
    wavefunction of the proton.}
  \label{fig:17.3}
\end{figure}

To derive a first approximation to the cross section for electron-proton
scattering, we consider this reaction from a frame in which the electron and
proton are moving rapidly toward each other, for example, the electron-proton
center-of-mass frame.  We assume that the center-of-mass energy is large
enough that we can ignore the proton mass in working out the kinematics.
Then the proton has an almost lightlike momentum along the collision axis.
The constituents of the proton also have lightlike momenta, which are almost
collinear with the momentum of the proton.  This is because a constituent
cannot acquire a large transverse momentum except through exchange of a hard
gluon, a process that is suppressed by the smallness of $\alpha_s$ at large
momentum scales.  Thus, to leading order in QCD perturbation theory, we can
write
\begin{equation}
  p = \xi P ,
  \label{eq:17.19}
\end{equation}
where $\xi$ is a number between 0 and 1, called the \emph{longitudinal
fraction} of the constituent.  To leading order in $\alpha_s$ we can also
ignore gluon emission or exchanges during the collision process.  The cross
section for electron-proton scattering is then given by the cross section for
electron-quark scattering at given longitudinal fraction $\xi$, multiplied by
the probability that the proton contains a quark at that value of $\xi$,
integrated over $\xi$.

But the probability that the proton contains a certain constituent with a
certain momentum fraction cannot be computed using QCD perturbation theory,
since it depends on the soft processes that determine the structure of the
proton as a bound state of quarks and gluons.  We will therefore consider
this probability to be an unknown function, to be determined from experiment.
Eventually, we will need to make use of such a probability function for each
species of quark, antiquark, and gluon that can be found in the wavefunction
of the proton.  Collectively, these constituents are called \emph{partons}.
For each parton species $f$, we write the probability of finding a
constituent of the proton of type $f$ at longitudinal fraction $\xi$ as
\begin{equation}
  f_f(\xi)\,d\xi = \text{probability of finding a parton of type } f
    \text{ with longitudinal fraction } \xi .
  \label{eq:17.20}
\end{equation}
The functions $f_f(\xi)$ are called the \emph{parton distribution functions}.
Using this notation, the cross section for electron-proton inelastic
scattering is given to leading order in $\alpha_s$ by the expression
\begin{equation}
  \sigma\bigl( e^{-}(k)p(P) \to e^{-}(k') + X \bigr)
  = \int_0^1 d\xi\ \sum_f f_f(\xi)\ \hat{\sigma}\bigl( e^{-}(k)q_f(\xi P)
    \to e^{-}(k') + q_f(p') \bigr) ,
  \label{eq:17.21}
\end{equation}
where $X$ stands for any hadronic final state.  The sum in \eqref{eq:17.21}
contains contributions from constituent antiquarks as well as constituent
quarks.

Equation \eqref{eq:17.21} is equivalent to the formula \eqref{eq:14.8} that
we constructed for this reaction in Chapter~14.  Now we see that this formula
is justified by the smallness of the QCD coupling constant at large momentum
scales.  It is important to remember, however, that \eqref{eq:17.21} is not
the complete prediction of QCD, but only the first term of an expansion in
$\alpha_s$; this level of approximation is called the \emph{parton model}.
The higher-order QCD corrections to Eq.~\eqref{eq:17.21} will involve
modifications both to the electron-quark scattering cross section and to the
parton distribution functions.  The most important of these corrections are
discussed in Section~17.5.

In the same way, all other reactions of the proton that involve large momentum
transfer also have parton model descriptions.  In QCD, all of these reaction
cross sections are computed from scattering amplitudes for quarks and gluons.
The initial motion of the partons for any process is described by the same
parton distribution functions $f_f(\xi)$ that appear in deep inelastic
scattering.

Let us now work out the explicit leading-order formula for the deep inelastic
scattering cross section, reviewing the analysis in Chapter~14.  In
Eq.~\eqref{eq:14.3}, we wrote the leading-order differential cross section
for the parton-level process,
\begin{equation}
  \frac{d\sigma}{dt}\bigl( e^{-}q \to e^{-}q \bigr)
  = \frac{2\pi\alpha^2 Q_f^2}{s^2}\ \frac{s^2 + u^2}{t^2} .
  \label{eq:17.22}
\end{equation}
In general, we will use the symbols $s$, $t$, $u$ to denote the Mandelstam
variables for two-body scattering processes at the parton level.  These
variables must be related to observable properties of the hadronic system or
the scattered electron.  For massless initial and final particles,
\begin{equation}
  s + t + u = 0 .
\end{equation}
In the case of deep inelastic scattering,
\begin{equation}
  t = -Q^2
  \qquad\text{and}\qquad
  s = 2p\cdot k = 2\xi P\cdot k = \xi s .
\end{equation}
Combining these relations with \eqref{eq:17.22} and the parton model
factorization \eqref{eq:17.21}, we obtain the deep inelastic scattering cross
section:
\begin{equation}
  \frac{d^2\sigma}{dx\,dQ^2}\bigl( e^{-}p \to e^{-}X \bigr)
  = \frac{2\pi\alpha^2}{Q^4}\ \Bigl[ 1 + (1-y)^2 \Bigr]\ \sum_f Q_f^2\,
    f_f(x) ,
  \label{eq:17.30}
\end{equation}
where $x = Q^2/2P\cdot q$ is the Bjorken scaling variable and
$y = 1 - Q^2/xs$.  This formula exhibits the property of Bjorken scaling, in
which the structure function depends only on $x$ and not on $Q^2$, up to
corrections of order $\alpha_s$.

The expression \eqref{eq:17.30} also predicts a specific dependence on $y$:
the factor $[1 + (1-y)^2]$ is characteristic of the scattering of electrons
from massless fermions.  This relation gave evidence that the partons involved
in deep inelastic scattering were fermions, at a time when the relation
between partons and quarks was still unclear.

\subsection{Deep Inelastic Neutrino Scattering}

Because the sum over quark flavors factorizes in Eq.~\eqref{eq:17.30}, one
cannot determine the individual parton distribution functions $f_f(x)$ from
electron scattering experiments alone.  One can, however, obtain more
detailed information on the structure of the proton through deep inelastic
neutrino scattering.

Neutrinos have zero electric charge and so do not interact by photon
exchange, but they do interact with quarks through weak interactions.  We
will discuss the weak interactions in detail in Chapter~20; for the moment,
let us adopt a simplified description that concentrates on the elementary
process shown in Fig.~\ref{fig:17.4}.  In this process, a neutrino converts
to the associated charged lepton,%
\footnote{There is also a neutral-current weak interaction in which the
neutrino remains a neutrino; see Problem 20.4.} exchanging a virtual massive
vector boson, the $W^{+}$.  This boson couples to a quark current that
converts a $d$ quark to a $u$ quark.  The effect of this exchange process is
to provide a different, but completely characterized, method for injecting a
hard momentum transfer $q$.  The amplitude for this process is described by
the effective Lagrangian
\begin{equation}
  \Delta\mathcal{L} = -\frac{g^2}{m_W^2}\ \bar{\ell}\gamma^\mu
    \frac{1-\gamma^5}{2}\nu\ \bar{u}\gamma_\mu\frac{1-\gamma^5}{2}d
    + \mathrm{h.c.},
  \label{eq:17.31}
\end{equation}
where $\ell$, $\nu$, $d$, $u$ are the fermion fields associated with the
charged lepton, the neutrino, and the $d$ and $u$ quarks, and $g$ is the weak
interaction coupling constant.  The factor $1/m_W^2$ comes from the $W$ boson
propagator, considered in the limit $q^2 \ll m_W^2$.  The first two factors
are often written in terms of the Fermi constant $G_F$, defined as
\begin{equation}
  \frac{G_F}{\sqrt{2}} = \frac{g^2}{8m_W^2} .
  \label{eq:17.32}
\end{equation}
This constant gives the strength of the weak interactions at energies much
less than $m_W$.  The crucial property of the weak interactions, shown
explicitly in \eqref{eq:17.31}, is that the $W$ boson couples only to the
\emph{left-handed} helicity states of relativistic fermions.  The deeper
significance of this property will be discussed in Chapter~20.

\begin{figure}[t]
  \centering
  \includegraphics[width=0.5\textwidth]{fig17.4.png}
  \caption{The elementary neutrino scattering process mediated by the weak
    interaction.}
  \label{fig:17.4}
\end{figure}

For technical reasons, it is easiest to do neutrino deep inelastic scattering
using muon neutrinos, which convert to muons after emitting a $W$ boson.  It
is equally feasible to scatter muon antineutrinos from nuclear targets, and,
as we will see, it is interesting to compare the effects of neutrinos and
antineutrinos.  Since the proton contains a small admixture of the heavier
quarks $s$, $c$, these also give small contributions to neutrino deep
inelastic scattering.  However, we will ignore these contributions in our
discussion.

The cross section for neutrino deep inelastic scattering is given by a
formula analogous to \eqref{eq:17.21}, with the photon-exchange cross
section replaced by one resulting from $W$ exchange.  It is straightforward
to work out this cross section directly.  However, we can also obtain the
result from Eq.~\eqref{eq:17.22}, if we look back to Chapter~5 and recall how
the structure of this equation arises from the various helicity
contributions.  In \eqref{eq:17.22}, the factor $t^2$ in the denominator came
from the photon propagator; this factor is replaced by $m_W^4$ in the weak
interaction case.  The factor $\tfrac{1}{2}s^2$ came from the Dirac matrix
algebra.  We saw in Section~5.2 that the first term is the contribution of
left-handed electrons scattering from left-handed fermions or right-handed
electrons scattering from right-handed fermions, and that the second term
arises from the other helicity combinations.  For the case of neutrino-quark
scattering, the interaction \eqref{eq:17.31} allows only the scattering of
left-handed neutrinos from left-handed quarks, so only the $s^2$ term appears.
To determine the overall normalization of the cross section, we note that,
since the neutrinos are produced in weak interactions, they always have
left-handed polarization, so no polarization average should be done.  On the
other hand, we must still average over the polarization of the initial quark.
In all, we find
\begin{equation}
  \frac{d\sigma}{dt}\bigl( \nu_\mu d \to \mu^{-}u \bigr)
  = \frac{G_F^2}{2(4\pi)^2}\ \frac{s^2}{t^2}\ \cdots
  = \frac{G_F^2}{2\pi}\,s\,(1-y)^2 .
  \label{eq:17.33}
\end{equation}
It is easy to check this formula by explicit computation starting from
\eqref{eq:17.31}.

The cross section for the scattering of antineutrinos from quarks can be
worked out in the same way.  Note that this reaction involves the exchange of
a $W^{-}$, and so converts $u$ quarks to $d$ quarks.  However, the $u$ quarks
must still be left-handed.  The only modification from the previous paragraph
comes in the fact that the antineutrinos that couple to the interaction
\eqref{eq:17.31} are right-handed, so the cross section comes from the term
in \eqref{eq:17.22} proportional to $u^2$:
\begin{equation}
  \frac{d\sigma}{dt}\bigl( \bar{\nu}_\mu u \to \mu^{+}d \bigr)
  = \frac{G_F^2}{2(4\pi)^2}\ \frac{u^2}{s^2}\ \cdots
  = \frac{G_F^2}{2\pi}\,s\,(1-y)^2 .
  \label{eq:17.34}
\end{equation}
Again, it is easy to verify this formula directly.  The cross section for
neutrino-antiquark scattering, converting a $\bar{u}$ into a $\bar{d}$, is
also given by Eq.~\eqref{eq:17.34}, while the cross section for
antineutrino-antiquark scattering, converting a $\bar{d}$ into a $\bar{u}$,
is given by Eq.~\eqref{eq:17.33}.

To convert these parton-level cross sections to physical cross sections, we
combine them with the parton distribution functions.  The kinematics is
exactly the same as in the case of electron scattering.  Thus, following the
arguments that led to Eq.~\eqref{eq:17.30}, we obtain the expressions
\begin{equation}
  \begin{aligned}
    \frac{d^2\sigma}{dx\,dy}\bigl( \nu p \to \mu^{-}X \bigr)
    &= \frac{G_F^2}{\pi}\,s\ \bigl[ xf_d(x) + xf_u(x)\cdot(1-y)^2 \bigr],\\
    \frac{d^2\sigma}{dx\,dy}\bigl( \bar{\nu} p \to \mu^{+}X \bigr)
    &= \frac{G_F^2}{\pi}\,s\ \bigl[ xf_u(x)\cdot(1-y)^2 + xf_d(x) \bigr].
  \end{aligned}
  \label{eq:17.35}
\end{equation}
According to these relations, deep inelastic neutrino scattering allows one to
map separately the parton distribution functions for $u$ and $d$ quarks and
antiquarks in the proton.

In addition, Eq.~\eqref{eq:17.35} makes a dramatic qualitative prediction: To
the extent that a proton (or neutron) is made of quarks with very few
additional quark-antiquark pairs, the deep inelastic neutrino scattering
cross section should be constant in $y$, while the antineutrino scattering
cross section should fall off as $(1-y)^2$.  The measured $y$ dependence of
these deep inelastic cross sections is shown in Fig.~\ref{fig:17.5}.  The
qualitative behavior predicted by the parton description is clearly evident;
the discrepancy from the strict prediction can be accounted for by a small
antiquark component in the nucleon wavefunction.

\begin{figure}[t]
  \centering
  \includegraphics[width=0.6\textwidth]{fig17.5.png}
  \caption{The distribution in $y$ of neutrino and anti-neutrino deep
    inelastic scattering from an iron target, as measured by the CDHS
    experiment, J.~G.~H.~de~Groot, et.\ al., \emph{Z.\ Phys.\ C} \textbf{1},
    143 (1979).  The solid curves are fits to the form
    $A + B(1-y)^2$.}
  \label{fig:17.5}
\end{figure}

\subsection{The Parton Distribution Functions}

Given that the parton model predictions for electron and neutrino deep
inelastic scattering do fit the data, one can make use of these relations to
extract the parton distribution functions and so learn something about the
structure of the proton.%
\footnote{A detailed discussion of the extraction of parton distribution
functions from data can be found in G.~Sterman, et.\ al., \emph{Rev.\ Mod.\
Phys.} \textbf{67}, 157 (1995).}  A set of distribution functions, chosen to
fit all available data, is shown in Fig.~\ref{fig:17.6}.  Since all of these
distributions, especially those for antiquarks, peak sharply at small $x$, we
have plotted $xf_f(x)$ for each species.  As we remarked in Chapter~14, a
small violation of Bjorken scaling is observed experimentally, so that these
distribution functions change slowly with $Q^2$.  The figure shows these
functions at $Q^2 = 4\ \mathrm{GeV}^2$.  We will see in Section~17.5 that
this violation of Bjorken scaling is an effect of higher-order corrections in
QCD; we will also argue that the measurement of this scaling violation allows
one to determine the parton distribution function for gluons, $f_g(x)$.
Anticipating this result, we have also plotted this function in the figure.
Not surprisingly, one finds that the $u$ and $d$ quarks are most likely to
carry a substantial fraction of the proton's momentum, while antiquarks and
gluons tend to have small longitudinal fractions.

\begin{figure}[t]
  \centering
  \includegraphics[width=0.6\textwidth]{fig17.6.png}
  \caption{Parton distribution functions $xf_f(x)$ for quarks, antiquarks,
    and gluons in the proton, at $Q^2 = 4\ \mathrm{GeV}^2$.  These
    distributions are obtained from a fit to deep inelastic scattering data
    performed by the CTEQ collaboration (CTEQ2L), described in J.~Botts,
    et.\ al., \emph{Phys.\ Lett.\ B} \textbf{304}, 159 (1993).}
  \label{fig:17.6}
\end{figure}

Since the parton distributions are the probabilities of finding various proton
constituents, they must be normalized in a way that reflects the quantum
numbers of the proton.  The proton is a bound state of $uud$, plus some
admixture of quark-antiquark pairs.  Thus it should contain an excess of two
$u$ quarks and one $d$ quark over the corresponding antiquarks.  These
considerations imply the constraints
\begin{equation}
  \int_0^1 dx\ \bigl[ f_u(x) - f_{\bar{u}}(x) \bigr] = 2,\qquad
  \int_0^1 dx\ \bigl[ f_d(x) - f_{\bar{d}}(x) \bigr] = 1 .
  \label{eq:17.36}
\end{equation}

So far we have discussed the parton distributions only for the proton.
Similar considerations, however, apply to any other hadron.  Each hadron has
its own set of parton distribution functions; these obey sum rules analogous
to \eqref{eq:17.36} but reflecting the particular quantum numbers of the
hadron.  The parton distribution functions should also reflect the symmetries
that link different hadrons.  For example, since the neutron can be
generated, to a few percent accuracy, by interchanging the role of $u$ and
$d$ quarks in the proton, its distribution functions obey
\begin{equation}
  f_u^n(x) = f_d(x),\qquad f_d^n(x) = f_u(x),\qquad
  f_s^n(x) = f_s(x),\qquad \text{etc.}
  \label{eq:17.37}
\end{equation}
In these equations, and henceforth, a distribution function without a special
label refers to the proton.  The parton distribution functions of the
antiproton are given by the exact relations
\begin{equation}
  f_u^{\bar{p}}(x) = f_{\bar{u}}(x),\qquad
  f_{\bar{u}}^{\bar{p}}(x) = f_u(x),\qquad \text{etc.}
  \label{eq:17.38}
\end{equation}

In any case, the total amount of momentum carried by the partons must sum to
the total momentum of the hadron.  This implies
\begin{equation}
  \int_0^1 dx\ x\ \bigl[ f_u(x) + f_d(x) + f_{\bar{u}}(x) + f_{\bar{d}}(x)
    + f_g(x) \bigr] = 1 .
  \label{eq:17.39}
\end{equation}
The distribution functions of quarks and antiquarks in the proton, as
extracted from deep inelastic scattering data, contribute only about half of
the total value required for this integral.  The remaining energy-momentum
must be carried by the gluons.

\section{Hard-Scattering Processes in Hadron Collisions}

If one collides hadrons with other hadrons at very high energy, most of the
collisions will involve only soft interactions of the consituent quarks and
gluons.  Such interactions cannot be treated using perturbative QCD, because
$\alpha_s$ is large when the momentum transfer is small.  In some collisions,
however, two quarks or gluons will exchange a large momentum $p_\perp$
perpendicular to the collision axis.  Then, as in deep inelastic scattering,
the elementary interaction takes place very rapidly compared to the internal
time scale of the hadron wavefunctions, so the lowest-order QCD prediction
should accurately describe the process.  Again, we should find a parton-model
formula that is built from a leading-order subprocess cross section,
integrated with parton distribution functions.  For the case of proton-proton
scattering, these functions will be the same ones that are measured in
lepton-proton deep inelastic scattering.

For example, if the hard parton-level process involves quark-antiquark
scattering into a final state $Y$, the leading-order QCD prediction takes the
form
\begin{equation}
  \sigma\bigl( p(P_1) + p(P_2) \to Y + X \bigr)
  = \int_0^1 dx_1 \int_0^1 dx_2\ \sum_f f_f(x_1)f_{\bar{f}}(x_2)\cdot
    \hat{\sigma}\bigl( q_f(x_1P) + \bar{q}_{\bar{f}}(x_2P) \to Y \bigr) ,
  \label{eq:17.40}
\end{equation}
where the sum runs over all species of quarks and antiquarks---$u$, $d$,
$\bar{u}$, $\bar{d}$, \ldots.  (Here again, $X$ denotes any hadronic final
state.)  The same formula, with appropriately modified distribution
functions, applies to any other hadron-hadron collision.  This formula will
be a good first approximation if, by some invariant measure, a large momentum
is transfered in the $q\bar{q}$ reaction.  In this section we will discuss
several examples of processes of this type.

\subsection{Lepton Pair Production}

The first and simplest example of a hadron-hadron hard-scattering process is
lepton pair production, also called the Drell-Yan process.  In this process, a
quark from one hadron and an antiquark from the other annihilate into a
virtual photon, which materializes as a pair of leptons:
\begin{equation}
  q\bar{q} \to \gamma^{*} \to \mu^{+}\mu^{-} .
\end{equation}
This is the hadronic analogue of the reaction $e^{+}e^{-}\to\mu^{+}\mu^{-}$
that we studied in Chapter~5.  The parton-level cross section for
$q\bar{q}\to\mu^{+}\mu^{-}$ is obtained from the $e^{+}e^{-}$ formula by
replacing the electron charge with the quark charge:
\begin{equation}
  \hat{\sigma}\bigl( q\bar{q} \to \mu^{+}\mu^{-} \bigr)
  = Q_f^2\,\frac{4\pi\alpha^2}{3\hat{s}} .
  \label{eq:17.41}
\end{equation}
Combining this with the parton model formula \eqref{eq:17.40}, we obtain the
Drell-Yan cross section.  The invariant mass of the lepton pair is
\begin{equation}
  M^2 = x_1x_2\,s ,
\end{equation}
and the longitudinal momentum of the pair determines the difference
$x_1 - x_2$.  The cross section for lepton pair production is therefore
\begin{equation}
  \frac{d\sigma}{dM^2}\bigl( pp \to \mu^{+}\mu^{-} + X \bigr)
  = \frac{4\pi\alpha^2}{9sM^2}\ \int dx_1\,dx_2\ \sum_f Q_f^2\,
    \bigl[ f_f(x_1)f_{\bar{f}}(x_2) + f_{\bar{f}}(x_1)f_f(x_2) \bigr]\,
    \delta(x_1x_2 - M^2/s) ,
  \label{eq:17.48}
\end{equation}
where $x_1$ and $x_2$ are given by the constraints of the kinematics.  It is
remarkable that the cross section for the Drell-Yan process is determined
point by point by information derivable from deep inelastic scattering.
Unfortunately, the relation between the two processes implied by
\eqref{eq:17.48} receives a correction of order $\alpha_s(M)$ that turns out
to be numerically large, and which must be included to check this prediction
against experiment.

\subsection{General Kinematics of Pair Production}

In deriving Eq.~\eqref{eq:17.48}, we used the total cross section
\eqref{eq:17.41} for the parton-level process, integrated over the angular
distribution of the outgoing leptons.  In principle, we could have retained
the angular information and derived a triply differential distribution.  This
would be the most complete prediction possible for a two-body parton-level
reaction.  It will be useful to work out the kinematics of such reactions,
taking a more general viewpoint.

In the generic situation, a parton of type 1 from proton 1 scatters from a
parton of type 2 from proton 2, yielding partons of types 3 and 4, with a
squared momentum transfer $\hat{t}$.  This generic process is shown in
Fig.~\ref{fig:17.8}.  In the Drell-Yan process, partons 3 and 4 are leptons.
But these partons could also be quarks or gluons, which materialize as
hadronic jets.  We assume that all partons can be treated as massless.  In
parton variables, the cross section for this process is
\begin{equation}
  \frac{d\sigma}{dx_1\,dx_2\,d\hat{t}}\bigl( p_1p_2 \to 3 + 4 + X \bigr)
  = f_1(x_1)f_2(x_2)\ \frac{d\hat{\sigma}}{d\hat{t}}\bigl( 1 + 2 \to 3 + 4
    \bigr) .
  \label{eq:17.49}
\end{equation}

\begin{figure}[t]
  \centering
  \includegraphics[width=0.5\textwidth]{fig17.8.png}
  \caption{A generic two-body parton scattering process.}
  \label{fig:17.8}
\end{figure}

Let us now translate this formula to observable parameters of the final
state.  In the leading order of QCD, the transverse momenta of partons 3 and
4 must be equal and opposite, but their longitudinal momenta are not
constrained.  We will take the three parameters of the final state to be the
common magnitude of the parton transverse momenta $p_\perp$ and the
longitudinal rapidities $y_3$, $y_4$ of the final-state partons, defined by
the formulae
\begin{equation}
  E_i = p_\perp\cosh y_i;\qquad
  p_{i\parallel} = p_\perp\sinh y_i .
  \label{eq:17.50}
\end{equation}
The longitudinal rapidity $y_i$ gives the boost of the particle $i$ from the
frame where it has zero longitudinal momentum.%
\footnote{In the literature on hadron collisions, $y_i$ is usually called
simply the rapidity, with the restriction to longitudinal boosts being
understood.}  Recall from Section~3.3 that rapidities simply add under
collinear boosts.  The transverse momentum is invariant under longitudinal
boosts.  Thus, $(y_3, y_4, p_\perp)$ is a set of variables with convenient
Lorentz transformation properties with respect to boosts along the collision
axis.  We will now see that these three parameters are related in a
straightforward way to the underlying variables $x_1$, $x_2$, $\hat{t}$.

Consider the center of mass frame of the colliding partons.  The total energy
in this frame is $\sqrt{\hat{s}}$.  Let us use a subscript $*$ to denote
other quantities measured in this frame, for instance, $\theta_*$ for the
parton scattering angle.  Then
\begin{equation}
  p_{3\parallel*} = \frac{1}{2}\sqrt{\hat{s}}\cos\theta_*,\qquad
  p_{3\perp*} = \frac{1}{2}\sqrt{\hat{s}}\sin\theta_* ,
  \label{eq:17.51}
\end{equation}
and $p_{4*}$ is oriented just oppositely.  This frame is also the center of
mass frame of partons 3 and 4, so
\begin{equation}
  y_{3*} = -y_{4*} = y_* .
  \label{eq:17.52}
\end{equation}
Since rapidities transform by shifts, we can solve for $y_*$ and for the
rapidity $Y$ by which one must boost to reach this frame:
\begin{equation}
  y_* = \frac{1}{2}(y_3 - y_4),\qquad
  Y = \frac{1}{2}(y_3 + y_4) .
  \label{eq:17.53}
\end{equation}
The scattering angle $\theta_*$ is determined from $y_*$ by combining
\eqref{eq:17.51} with the relation $E_* = p_\perp\cosh y_*$:
\begin{equation}
  \frac{p_\perp}{\frac{1}{2}\sqrt{\hat{s}}} = \frac{1}{\sin\theta_*}
  = \cosh y_* .
  \label{eq:17.54}
\end{equation}
Then the Mandelstam variables
\begin{equation}
  \hat{s} = \frac{4p_\perp^2}{\sin^2\theta_*},\qquad
  \hat{t} = -\frac{1}{2}\hat{s}(1-\cos\theta_*)
  \label{eq:17.55}
\end{equation}
can be expressed as
\begin{equation}
  \hat{s} = 4p_\perp^2\cosh^2 y_*,\qquad
  \hat{t} = -2p_\perp^2\cosh y_*\,e^{-y_*} .
  \label{eq:17.56}
\end{equation}
We can combine the first of these expressions with $\hat{s} = x_1x_2s$ to
determine $x_1$ and $x_2$:
\begin{equation}
  x_1 = \frac{2p_\perp}{\sqrt{s}}\,\cosh y_*\,e^{Y},\qquad
  x_2 = \frac{2p_\perp}{\sqrt{s}}\,\cosh y_*\,e^{-Y} .
  \label{eq:17.57}
\end{equation}
To translate the cross section \eqref{eq:17.49} to the final parton
observables, we need the Jacobian
\begin{equation}
  \frac{\partial(x_1,x_2,\hat{t})}{\partial(y_3,y_4,p_\perp^2)} .
  \label{eq:17.58}
\end{equation}
Multiplying Eq.~\eqref{eq:17.49} by this factor gives
\begin{equation}
  \frac{d\sigma}{dy_3\,dy_4\,dp_\perp^2}\bigl( p_1p_2 \to 3+4+X \bigr)
  = f_1(x_1)f_2(x_2)\ \frac{d\hat{\sigma}}{d\hat{t}}\bigl( 1+2 \to 3+4
    \bigr)\ \cdots .
  \label{eq:17.59}
\end{equation}
This can be simplified a bit using the relations $\hat{s} = x_1x_2s$ and
$p_\perp dp_\perp = d^2p_\perp/2\pi$, yielding the final result:
\begin{equation}
  \frac{d\sigma}{dy_3\,dy_4\,d^2p_\perp}\bigl( p_1p_2 \to 3+4+X \bigr)
  = f_1(x_1)f_2(x_2)\ \frac{d\hat{\sigma}}{d\hat{t}}\bigl( 1+2 \to 3+4
    \bigr)\ \frac{1}{\hat{s}}\ \cdots .
  \label{eq:17.60}
\end{equation}
In this formula, $x_1$, $x_2$, and the Mandelstam variables of the parton
subprocess are given by Eqs.~\eqref{eq:17.57} and \eqref{eq:17.56}.

This result gives us the complete distribution of final-state leptons or jets
in any two-body reaction of partons.  For example, to find the distribution
of final-state leptons in the Drell-Yan process, we would insert into this
formula the differential cross section for quark annihilation into leptons,
\begin{equation}
  \frac{d\hat{\sigma}}{d\hat{t}}\bigl( q\bar{q} \to \ell^{+}\ell^{-} \bigr)
  = \frac{2\pi\alpha^2 Q_f^2}{3\hat{s}^2}\ \frac{\hat{t}^2 + \hat{u}^2}
    {\hat{t}^2} .
  \label{eq:17.61}
\end{equation}
The formula can be applied equally well to other two-body parton reactions,
if we know the relevant parton-level differential cross sections.

\subsection{Jet Pair Production}

The most common two-body parton reactions are those of QCD, involving
quarks, gluons, or both.  Unfortunately, it is very difficult to distinguish
hadronic jets initiated by gluons from those initiated by quarks.  It is even
more difficult to determine experimentally whether the initial partons in a
hard-scattering process were quarks or gluons.  Thus, the predictions of QCD
for hard-scattering processes are most often quoted as cross sections for jet
production in hadronic collisions, summing over all possible reactions of
quarks, antiquarks, and gluons.  In any event, to derive these predictions,
we must work out the basic parton-parton cross sections.

The simple two-body scattering processes of quarks, antiquarks, and gluons
are the elementary processes of QCD perturbation theory, in the same sense
that the reactions studied in Chapter~5 are the elementary processes of QED
perturbation theory.  They are the basic hadronic hard-scattering reactions
that appear in QCD at the leading order in $\alpha_s$.  In the remainder of
this section, we will write down the cross section formulae for the various
possible quark and gluon subprocesses.  All of these cross sections will be of
order $\alpha_s^2$.  In practice, this $\alpha_s$ should be evaluated at a
typical momentum transfer of the reaction, for example, at $Q^2 = \hat{t}$.

\begin{figure}[t]
  \centering
  \includegraphics[width=0.45\textwidth]{fig17.9.png}
  \caption{Feynman diagram contributing to $ud\to ud$.}
  \label{fig:17.9}
\end{figure}

The simplest subprocess is the scattering of different species of quarks, for
example, $u + d \to u + d$.  At order $\alpha_s^2$, this process occurs via
the Feynman diagram shown in Fig.~\ref{fig:17.9}.  This process is analogous
to electron-muon scattering in QED, for which we wrote the cross section in
Eq.~\eqref{eq:17.22}:
\begin{equation}
  \frac{d\hat{\sigma}}{d\hat{t}}\bigl( e\mu \to e\mu \bigr)
  = \frac{2\pi\alpha^2}{\hat{s}^2}\ \frac{\hat{s}^2 + \hat{u}^2}{\hat{t}^2}
  .
  \label{eq:17.62}
\end{equation}
To convert this to the cross section for quark scattering in QCD, we need
only replace the QED coupling $e^2$ by $g^2$ times an $SU(3)$ group theory
factor.  The QCD diagram contains the factor
\begin{equation}
  t^a_{ii'}\,t^a_{jj'} ,
\end{equation}
where $i$, $i'$ are the initial and final colors of the $u$ quark and $j$,
$j'$ are the initial and final colors of the $d$ quark.  To compute the cross
section, we must square this factor, sum over final colors, and average over
initial colors.  This gives the factor
\begin{equation}
  \frac{1}{9}\ \mathrm{tr}\bigl[ t^at^b \bigr]\mathrm{tr}\bigl[ t^at^b \bigr]
  = \frac{1}{9}\ \bigl[ C(r) \bigr]^2\,\delta^{ab}\delta^{ab}
  = \frac{1}{9}\cdot\frac{1}{4}\cdot 8
  = \frac{2}{9} ,
  \label{eq:17.63}
\end{equation}
where we have used Eq.~\eqref{eq:15.78} and $C(r) = 1/2$ for the fundamental
representation of $SU(3)$.  Thus for $ud$ scattering,
\begin{equation}
  \frac{d\hat{\sigma}}{d\hat{t}}\bigl( ud \to ud \bigr)
  = \frac{4\pi\alpha_s^2}{9\hat{s}^2}\ \frac{\hat{s}^2 + \hat{u}^2}
    {\hat{t}^2} .
  \label{eq:17.64}
\end{equation}
The same formula applies for the scattering of any two different quarks, or,
by crossing, to the scattering of a quark and an antiquark of a different
species.  Crossing from the $t$ to the $s$ channel gives the cross section
for $q\bar{q}$ annihilation into a different species:
\begin{equation}
  \frac{d\hat{\sigma}}{d\hat{t}}\bigl( u\bar{u} \to d\bar{d} \bigr)
  = \frac{4\pi\alpha_s^2}{9\hat{s}^2}\ \frac{\hat{t}^2 + \hat{u}^2}
    {\hat{s}^2} .
  \label{eq:17.65}
\end{equation}

\begin{figure}[t]
  \centering
  \includegraphics[width=0.5\textwidth]{fig17.10.png}
  \caption{Feynman diagrams contributing to $u\bar{u}\to u\bar{u}$.}
  \label{fig:17.10}
\end{figure}

The scattering of a quark with an antiquark of the same species is more
complicated, since now there are two Feynman diagrams, shown in
Fig.~\ref{fig:17.10}, which interfere with one another.  The analogous QED
process is Bhabha scattering, $e^{+}e^{-}\to e^{+}e^{-}$, for which we worked
out the cross section in Problem~5.2:
\begin{equation}
  \frac{d\sigma}{dt}\bigl( e^{+}e^{-} \to e^{+}e^{-} \bigr)
  = \frac{2\pi\alpha^2}{s^2}\ \Bigl( \frac{s^2 + u^2}{t^2}
    + \frac{s^2 + t^2}{u^2} + \frac{4s^2}{tu} \Bigr) .
  \label{eq:17.66}
\end{equation}
However, it is not quite straightforward to transcribe this to QCD, because
different terms receive different color factors.  The color factors for the
$t$-channel, $u$-channel, and interference terms are different, and must be
computed separately.  Carrying out this computation, one finds
\begin{equation}
  \frac{d\hat{\sigma}}{d\hat{t}}\bigl( u\bar{u} \to u\bar{u} \bigr)
  = \frac{4\pi\alpha_s^2}{9\hat{s}^2}\ \Bigl( \frac{\hat{s}^2 + \hat{u}^2}
    {\hat{t}^2} + \frac{\hat{s}^2 + \hat{t}^2}{\hat{u}^2}
    - \frac{2}{3}\frac{\hat{s}^2}{\hat{t}\hat{u}} \Bigr) .
  \label{eq:17.67}
\end{equation}

The remaining parton subprocesses involve gluons.  The cross sections for
quark-gluon scattering and gluon-gluon scattering are computed in
Problem~17.3.  The results are:
\begin{equation}
  \frac{d\hat{\sigma}}{d\hat{t}}\bigl( qg \to qg \bigr)
  = \frac{4\pi\alpha_s^2}{9\hat{s}^2}\ \Bigl( \frac{\hat{s}^2 + \hat{u}^2}
    {\hat{t}^2} - \frac{2}{3}\frac{\hat{s}}{\hat{u}} - \frac{2}{3}
      \frac{\hat{u}}{\hat{s}} + \frac{\hat{s}^2 + \hat{t}^2}{\hat{u}^2}
    \Bigr) ,
  \label{eq:17.68}
\end{equation}
\begin{equation}
  \frac{d\hat{\sigma}}{d\hat{t}}\bigl( gg \to gg \bigr)
  = \frac{4\pi\alpha_s^2}{\hat{s}^2}\ \frac{3}{2}\ \Bigl( \frac{\hat{s}^2
    + \hat{u}^2}{\hat{t}^2} + \frac{\hat{s}^2 + \hat{t}^2}{\hat{u}^2}
    - \frac{4}{3}\frac{\hat{s}^2}{\hat{t}\hat{u}} \Bigr) .
  \label{eq:17.69}
\end{equation}

Using these parton-level cross sections, together with the parton model
formula \eqref{eq:17.40}, we can compute the cross section for two-jet
production in hadron collisions.  Fig.~\ref{fig:17.13} compares the leading
order QCD prediction for two-jet production with the invariant mass
distribution of two-jet events observed in high-energy $p\bar{p}$ collisions.

\begin{figure}[t]
  \centering
  \includegraphics[width=0.6\textwidth]{fig17.13.png}
  \caption{Two-jet invariant mass distribution in $p\bar{p}$ collisions at
    $E_{cm} = 1.8$ TeV, as measured by the CDF collaboration, F.~Abe,
    et.\ al., \emph{Phys.\ Rev.\ D} \textbf{48}, 998 (1993).  The measurement
    is compared to a leading-order QCD calculation using the CTEQ structure
    functions described in Fig.~\ref{fig:17.6}.  The three lower curves show
    the invariant mass distributions for the three components of the
    theoretical prediction: quark-quark (and antiquark) scattering,
    quark-gluon scattering, and gluon-gluon scattering.}
  \label{fig:17.13}
\end{figure}

The overall normalization of the theoretical prediction is uncertain by about
a factor of 2 due to the ambiguity of the choice of $Q^2$ used to evaluate
$\alpha_s(Q^2)$ in the parton cross sections, and due to similar ambiguities
in deriving parton distributions from deep inelastic scattering cross
sections.  This uncertainty is reduced to about 30\% when corrections of
order $\alpha_s$ are included.  Still, it is remarkable that the lowest-order
QCD prediction tracks the observed distribution as a function of the two-jet
invariant mass as it falls by six orders of magnitude.  Thus, for the jet
production cross section, as for hard processes involving leptons, QCD indeed
gives a reasonable description of the behavior of the strong interactions at
large momentum transfer.

\section{Parton Evolution}

Now that we have examined the predictions of QCD at the leading order for
several strong interaction processes, we should investigate the corrections to
these predictions at the next order in $\alpha_s$.  We saw in Section~17.2
that the corrections from individual diagrams may contain mass singularities,
singularities associated with collinear emission processes which appear in the
limit of zero mass.  For the process of $e^{+}e^{-}$ annihilation to hadrons,
we saw that these mass singularities, and the infrared divergences from soft
gluon emission, cancel in the expression for the \emph{total} cross section.
It can be shown that this is a general feature of processes in which quarks
and gluons are produced in the collision of leptons or photons.  However,
when quarks or gluons appear in the \emph{initial} state of a parton
subprocess, the corrections to the process will, in general, have mass
singularities that do not cancel.  In this section we will demonstrate this
effect and work out its physical interpretation.  We will find that these
singular terms predict a violation of Bjorken scaling by terms depending on
the logarithm of the momentum scale.  In fact, they lead to a precise set of
differential equations that govern the momentum dependence of the parton
distributions.

The basic phenomena associated with mass singularities in QCD are already
present in the physics of collinear photon emission in QED at high energies,
and so it is most straightforward to begin by studying that case.  In this
section, we will show that collinear photon emission leads to an analogue of
a parton distribution function for the electron.  We will derive a
differential equation describing this distribution function, first constructed
by Gribov and Lipatov.  Finally, we will generalize this equation to QCD,
following the construction of Altarelli and Parisi.%
\footnote{V.~N.~Gribov and L.~Lipatov, \emph{Sov.\ J.\ Nucl.\ Phys.}
\textbf{15}, 438 (1972); G.~Altarelli and G.~Parisi, \emph{Nucl.\ Phys.\ B}
\textbf{126}, 298 (1977).  We also strongly recommend reading the papers of
J.~Kogut and L.~Susskind, \emph{Phys.\ Rev.\ D} \textbf{9}, 697, 3391
(1974).}

In Chapters~5 and~6, we studied several examples of QED processes that
involved diagrams with $t$- or $u$-channel singularities.  In these cases, we
found that the total cross section was generally enhanced by an extra factor
$\log(s/m^2)$ in the high-energy limit.  For example, in Eq.~\eqref{eq:5.95}
we saw that the $u$-channel exchange diagram in Compton scattering,
Fig.~\ref{fig:17.14}(a), leads to an integral that, in the high-energy limit,
takes the form
\begin{equation}
  \int d\cos\theta\ \frac{1}{(1+\cos\theta)} .
\end{equation}
The singularity as $\cos\theta\to -1$ is cut off by the electron mass, leading
to the logarithmic enhancement factor.  Thus the collinear photon emission
costs a factor that is not $\alpha$ but rather $\alpha\log(s/m^2)$.  Emission
of multiple collinear photons, as in Fig.~\ref{fig:17.14}(b), gives
contributions of order $(\alpha\log(s/m^2))^n$.  To improve the accuracy of
perturbation theory, it would be useful to find a procedure for summing these
terms to all orders in $\alpha$.  In QCD, the corresponding factor for
collinear gluon emission would be
\begin{equation}
  \alpha_s(Q^2)\log\frac{Q^2}{\mu^2} ,
\end{equation}
where $\mu$ is the momentum scale where nonperturbative QCD effects become
important.  Comparing with Eq.~\eqref{eq:17.17}, we see that this product is
of order 1.  Thus, in this case, the resummation of large logarithms is
essential if we are to make any quantitative predictions.

\begin{figure}[t]
  \centering
  \includegraphics[width=0.55\textwidth]{fig17.14.png}
  \caption{Diagrams with mass singularities associated with collinear photon
    emission: (a) leading order; (b) higher order.}
  \label{fig:17.14}
\end{figure}

\begin{figure}[t]
  \centering
  \includegraphics[width=0.55\textwidth]{fig17.15.png}
  \caption{General form of diagrams with mass singularities in QED.}
  \label{fig:17.15}
\end{figure}

In QED, diagrams with mass singularities associated with one collinear
emission are of one of the forms shown in Fig.~\ref{fig:17.15}.  In each case,
the circle represents a scattering process with large momentum transfer.  The
mass singularity appears when the denominator of the intermediate propagator
vanishes, that is, when the intermediate state is almost on-shell.  Thus, it
is natural to consider the first diagram in Fig.~\ref{fig:17.15} to be a
transition to a real photon and an almost-real electron, followed by the
interaction of the electron with the remaining particles in the amplitude.
The second diagram should have a similar interpretation with an almost-real
photon in the intermediate state.  The only subtlety comes in defining the
polarization of the intermediate-state particle.  For the case of an
intermediate-state electron, the numerator of the propagator is
\begin{equation}
  \slashed{k} + m .
  \label{eq:17.79}
\end{equation}
Thus, when $k^2\to 0$, the photon emission vertex and the remaining part of
the amplitude are contracted with on-shell polarization spinors for a massless
electron.  The analogous statement for the diagram with the photon in the
intermediate state would be that the electron emission vertex and the
remaining photon amplitude should be contracted with physical transverse
polarization vectors for the intermediate-state photon.  Since the numerator
of the photon propagator is $g^{\mu\nu}$, it is not obvious that the photon
propagator reduces in this way.  But it is true.  To see this, use the
expansion for $g^{\mu\nu}$ in terms of massless polarization vectors given in
Eq.~(16.20):
\begin{equation}
  g^{\mu\nu} = \epsilon^{\mu}_{+}\epsilon^{\nu*}_{-} + \epsilon^{\mu}_{-}
    \epsilon^{\nu*}_{+} - g^{\mu\nu}\ \cdots .
  \label{eq:17.80}
\end{equation}
Here $\epsilon^{\mu}_{\pm}$ are transverse polarization vectors.  The forward
polarization vector $\epsilon_{+}$ is proportional to the photon momentum
$q^\mu$.  When we contract $\epsilon_{+}$ with the QED scattering amplitude on
the right, we will obtain zero by the Ward identity, and the contraction of
$\epsilon_{+}$ with the electron emission vertex similarly gives zero.  Thus,
for the purpose of computing the singular term as the photon momentum $q$ goes
on-shell, we may replace
\begin{equation}
  \frac{-ig^{\mu\nu}}{q^2} \to \frac{i}{q^2}\ \sum_{\lambda=\pm}
    \epsilon^{\mu*}_{\lambda}\epsilon^{\nu}_{\lambda} ,
  \label{eq:17.81}
\end{equation}
and evaluate the photon emission and absorption amplitudes with transverse
polarization vectors.

\subsection{Matrix Element for Electron Splitting}

By replacing the numerator of the intermediate propagator with a sum over
polarization vectors, we decouple the photon or electron emission vertex from
the rest of the diagram.  We will now evaluate this vertex explicitly between
physical polarization states of massless particles.  The kinematics is shown
in Fig.~17.16.  The two final particles should be almost collinear,
with a small relative transverse momentum.  We can choose the incident
electron momentum to lie along the 3 axis and the outgoing momenta to lie in
the 1-3 plane.  Let $z$ be the fraction of the energy of the initial electron
that is carried off by the photon.  Then the three 4-momenta can be written as
\begin{equation}
  \begin{aligned}
    p &= (p, 0, 0, p),\\
    q &\approx (zp, p_\perp, 0, zp),\\
    k &\approx ((1-z)p, -p_\perp, 0, (1-z)p) .
  \end{aligned}
  \label{eq:17.82}
\end{equation}
These three vectors satisfy $p^2 = q^2 = k^2 = 0$, up to terms of order
$p_\perp^2$.  In the process where a real photon is emitted, we should have
$p^2$ and $q^2$ exactly zero, and $k^2$ slightly off-shell by an amount of
order $p_\perp^2$.  We will need to know the value of $k^2$, which appears in
the virtual electron propagator.  So let us modify Eqs.~\eqref{eq:17.82} to
satisfy the condition $q^2 = 0$ up to terms of order $p_\perp^2$, rewriting
$q$ and $k$ as
\begin{equation}
  \begin{aligned}
    q &= \Bigl( zp, p_\perp, 0, zp - \frac{p_\perp^2}{2p} \Bigr),\\
    k &= \Bigl( (1-z)p, -p_\perp, 0, (1-z)p + \frac{p_\perp^2}{2p}
      \Bigr) .
  \end{aligned}
  \label{eq:17.83}
\end{equation}
With this modification,
\begin{equation}
  k^2 = -\frac{p_\perp^2}{1-z}\ -\ 2(1-z)\,\cdots = -\frac{p_\perp^2}{1-z}
    + \mathcal{O}(p_\perp^4) .
  \label{eq:17.84}
\end{equation}
Thus, if the photon is real and the electron is virtual, we have
\begin{equation}
  q^2 = 0,\qquad k^2 = -\frac{p_\perp^2}{1-z} .
\end{equation}
Reciprocally, in the process with a real electron and a virtual photon,
\begin{equation}
  k^2 = 0,\qquad q^2 = -\frac{p_\perp^2}{z} .
  \label{eq:17.85}
\end{equation}
These more accurate expressions will be needed only in the propagator of the
virtual particle.  The matrix element of the electron-photon vertex begins in
order $p_\perp$, so it is not significantly affected by the modification of
\eqref{eq:17.82} to \eqref{eq:17.83}, and is the same (to lowest order) no
matter which particle is virtual.

We now calculate the matrix elements of the QED vertex between massless
states of definite helicity.  If the initial electron is left-handed, the
final electron must also be left-handed, by helicity conservation.  Then the
photon emission vertex is given by
\begin{equation}
  i\mathcal{M} = \bar{u}_L(k)\,(-ie\gamma^\mu)\,u_L(p)\,
    \epsilon^{*}_\mu(q) ,
  \label{eq:17.86}
\end{equation}
where the photon polarization vector may be either left- or right-handed.
Recalling the helicity-basis expressions
\begin{equation}
  u_L(p) = \begin{pmatrix}\sqrt{p\cdot\sigma}\,\xi_L\\ \sqrt{p\cdot
    \bar{\sigma}}\,\xi_L\end{pmatrix}
  \qquad\text{for } m = 0 ,
\end{equation}
we can write more explicitly
\begin{equation}
  i\mathcal{M} = ie\sqrt{2(1-z)p}\,\sqrt{2p}\ \xi^{\dagger}_L(k)\,
    \bm{\sigma}\cdot\bm{\epsilon}^{*}(q)\,\xi_L(p) .
  \label{eq:17.87}
\end{equation}
Evaluating the spinor products, we find that the amplitude is nonzero only
when the photon has the same helicity as the electron (for the case of a
left-handed electron emitting a photon, the photon must be right-handed, by
angular momentum conservation).  The result, after squaring and summing over
polarizations, gives the splitting probability
\begin{equation}
  P_{e\to e\gamma}(z) = \frac{\alpha}{2\pi}\ \frac{1+z^2}{1-z} .
\end{equation}

Following the steps that led to Eq.~(17.97), we find that the cross
section for the process in which the electron emits a collinear photon and
then scatters is
\begin{equation}
  \begin{aligned}
    \sigma(e^{-}\to Y) &= \int_0^1 dz\ \int \frac{dp_\perp^2}{p_\perp^2}\
      \frac{\alpha}{2\pi}\ \frac{1+(1-z)^2}{z}\ \hat{\sigma}(e^{-} \to Y)
      \ \cdots\\
    &= \int dz\ \int \frac{dp_\perp^2}{p_\perp^2}\ \frac{\alpha}{2\pi}\
      \frac{1+(1-z)^2}{z}\ \hat{\sigma}(e^{-} \to Y) ,
  \end{aligned}
  \label{eq:17.99}
\end{equation}
where the intermediate electron carries a longitudinal fraction $(1-z)$.

It is tempting to substitute $x = (1-z)$ and interpret the factor multiplying
the cross section under the integral in \eqref{eq:17.99} as the parton
distribution for finding an electron parton in the electron.  This would give
\begin{equation}
  f^{(1)}_e(x) = \frac{\alpha}{2\pi}\log\frac{Q^2}{m^2}\ \frac{1+x^2}{1-x}
  .
  \label{eq:17.100}
\end{equation}
However, this expression is not adequate.  Most obviously, it does not take
into account the processes without radiation, in which the electron remains
an electron at the full energy.  This is easily remedied by considering
\eqref{eq:17.100} as the order-$\alpha$ correction to the most naive
expectation,
\begin{equation}
  f^{(0)}_e(x) = \delta(1-x) ,
  \label{eq:17.101}
\end{equation}
in which we consider the electron to contain only an electron at the full
energy.  Unfortunately, the sum of \eqref{eq:17.101} and \eqref{eq:17.100}
still does not give an adequate description of the electron distribution, for
two reasons.  First, Eq.~\eqref{eq:17.100} diverges near $x = 1$, and we need
a prescription for treating this singularity.  Second, while Eq.~\eqref{eq:17.100}
takes into account the virtual electrons moved to longitudinal fraction $x$
from $x = 1$ by radiation, it does not take into account the concomitant loss
of electrons from the delta function peak at $x = 1$.

The divergence of \eqref{eq:17.100} at $x = 1$ corresponds to the emission of
soft photons.  We saw in Section~6.5 that the emission of soft photons does
not affect the rate of a QED reaction.  Order by order in $\alpha$, one finds
that infrared-divergent positive contributions to the total rate from the
emission of soft photons are balanced by negative contributions from diagrams
with soft virtual photons.  In the present example, the negative contribution
must decrease the weight of the process in which no photon is emitted.
Thus, to order $\alpha$, the parton distribution for electrons in the electron
should have the form
\begin{equation}
  f_e(x) = \delta(1-x) + \frac{\alpha}{2\pi}\log\frac{Q^2}{m^2}\ \Bigl[
    \frac{1+x^2}{(1-x)_+} - A\,\delta(1-x) \Bigr] .
  \label{eq:17.102}
\end{equation}
The coefficient $A$ results from diagrams with virtual photons that we have
not computed.  However, the effect of these diagrams is easy to understand;
they subtract from the delta function the probability that has been moved to
lower $x$ by radiation, so that the integral over the full term of order
$\alpha$ is zero.  Another way of expressing this criterion is that $A$ is
determined by the condition that the electron contain exactly one electron
parton,
\begin{equation}
  \int_0^1 dx\ f_e(x) = 1 .
  \label{eq:17.103}
\end{equation}
(This equation will be modified below, when we include pair-creation
processes.)

It is not so clear how to integrate over the singular denominator in
\eqref{eq:17.100} to determine $A$ explicitly.  It is conventional to define a
distribution that can be integrated by subtracting a delta function from the
singular term.  Define the distribution
\begin{equation}
  \frac{1}{(1-x)_+}
  \label{eq:17.104}
\end{equation}
to agree with the function $1/(1-x)$ for all values of $x$ less than 1, and
to have a singularity at $x = 1$ such that the integral of this distribution
with any smooth function $f(x)$ gives
\begin{equation}
  \int_0^1 dx\ \frac{f(x)}{(1-x)_+} = \int_0^1 dx\ \frac{f(x) - f(1)}
    {1-x} .
  \label{eq:17.105}
\end{equation}
Less formally,
\begin{equation}
  \frac{1}{(1-x)_+} = \lim_{\epsilon\to 0}\ \Bigl[ \frac{1}{1-x}
    \theta(1-\epsilon-x) - \delta(1-x)\ \int_0^{1-\epsilon} \frac{d\xi}{1-\xi}
    \Bigr] .
  \label{eq:17.106}
\end{equation}
The more formal definition \eqref{eq:17.105} is often easier to use in
practice.

Using this definition, we can bring a piece of the delta function into the
singular term of \eqref{eq:17.102} by changing the denominator $(1-x)$ to
$(1-x)_+$.  Then, to normalize \eqref{eq:17.102}, we need the integral
\begin{equation}
  \int_0^1 dx\ \frac{1+x^2}{1-x} = \int_0^1 dx\ \Bigl( \frac{x^2-1}{1-x}
    + \frac{2}{1-x} \Bigr) .
\end{equation}
Our final form of the electron distribution, to order $\alpha$, is
\begin{equation}
  f_e(x) = \delta(1-x) + \frac{\alpha}{2\pi}\log\frac{Q^2}{m^2}\ \Bigl[
    \frac{1+x^2}{(1-x)_+} - \frac{3}{2}\delta(1-x) \Bigr] .
  \label{eq:17.107}
\end{equation}
This distribution is now properly normalized, but it is still highly singular
near $x = 1$.  Thus, we should expect higher-order corrections to the
electron distribution function to be important in this region.  We must,
then, think about how to treat the emission of many collinear photons.

\subsection{Multiple Splittings}

In fact, it is not difficult to extend the analysis we have just completed to
account for emission of many collinear photons.  Consider the process shown
in Fig.~\ref{fig:17.17}(a), in which photon 1 is radiated with a transverse
momentum $p_{1\perp}$ and photon 2 with a transverse momentum $p_{2\perp}$.
The emission of photon 2 can be computed just as we did above.  If
$p_{2\perp} \gg p_{1\perp}$, the first virtual electron is very close to mass
shell compared to $p_{2\perp}$ and so we can ignore its virtuality in
computing the emission of the photon 1.  The double photon emission gives a
contribution of order
\begin{equation}
  \alpha^2\ \log^2\frac{Q^2}{m^2} .
\end{equation}
In the opposite limit, $p_{2\perp} \ll p_{1\perp}$, there is no denominator of
order $p_{2\perp}^2$, and so we do not find a double logarithm.  Only in the
case $p_{2\perp} \ll p_{1\perp}$ can the contribution of order $\alpha^2$
compete with the contribution of order $\alpha$.

\begin{figure}[t]
  \centering
  \includegraphics[width=0.55\textwidth]{fig17.17.png}
  \caption{Higher-order diagrams with collinear photon emission: (a) two
    collinear photons; (b) many collinear photons.}
  \label{fig:17.17}
\end{figure}

This argument extends to the emission of arbitrarily many collinear photons,
Fig.~\ref{fig:17.17}(b).  The region of integration over the photon phase
space corresponding to the ordering
\begin{equation}
  p_{1\perp} \gg p_{2\perp} \gg p_{3\perp} \gg \cdots
  \label{eq:17.108}
\end{equation}
gives a contribution that contains the factor
\begin{equation}
  \frac{1}{n!}\ \alpha^n\ \log^n\frac{Q^2}{m^2} .
  \label{eq:17.109}
\end{equation}
If the photon transverse momenta are ordered in any other way, the
contribution from that region contains at least one less power of the large
logarithm at the same order in $\alpha$.  If condition \eqref{eq:17.108}
holds, the virtual electron momenta are increasingly off-mass-shell as one
proceeds from the outside of the diagram toward the hard collision.  In this
case, the electron momenta are said to be \emph{strongly ordered}.

This set of conclusions has an interesting physical interpretation.  Since the
intermediate electrons are increasingly virtual as we go into the diagram, it
is natural to interpret them as components of the physical electron when this
particle is analyzed at successively smaller distance scales.  The
intermediate electron with $k^2 \sim p_\perp^2$ may be thought of as a
constituent of the electron made visible when the wavefunction of the
physical electron is probed with a resolution $\Delta r \sim (p_\perp)^{-1}$.
In this picture, the electron seen at one resolution can be resolved at a
finer scale into a more virtual electron and a number of photons.

From either the perspective of computing Feynman diagrams or the grander
perspective of the electron structure, it is useful to imagine the splitting
of the electron into a virtual electron plus photons to be a continuous
evolution process as a function of the transverse momentum of the electron
constituent.  To describe this process mathematically, we introduce an
explicit $p_\perp$ dependence of the electron and photon distribution
functions.  We define the functions $f_\gamma(x, Q)$ and $f_e(x, Q)$ to give
the probabilities of finding a photon or an electron of longitudinal fraction
$x$ in the physical electron, taking into account collinear photon emissions
with transverse momenta $p_\perp < Q$.

If $Q$ is slightly increased to $Q + \Delta Q$, we must take into account the
possibility that an electron constituent in $f_e(x, Q)$ will radiate a photon
with $Q < p_\perp < Q + \Delta Q$.  The differential probability for an
electron to split off a photon that carries away a fraction $z$ of its energy
is
\begin{equation}
  \frac{\alpha}{2\pi}\ \frac{dp_\perp^2}{p_\perp^2}\ \frac{1+(1-z)^2}{z} .
  \label{eq:17.110}
\end{equation}
The new photon distribution can therefore be computed as follows:
\begin{equation}
  \begin{aligned}
    f_\gamma(x, Q + \Delta Q) &= f_\gamma(x, Q) + \frac{\alpha}{2\pi}\
      \frac{\Delta Q^2}{Q^2}\ \int_0^1 dx'\,dz\ \frac{1+(1-z)^2}{z}\,
      f_e(x', p_\perp)\,\delta(x - zx')\\
    &= f_\gamma(x, Q) + \frac{\alpha}{2\pi}\ \frac{\Delta Q^2}{Q^2}\
      \int_x^1 \frac{dz}{z}\ \frac{1+(1-z)^2}{z}\ f_e\Bigl(
        \frac{x}{z}, Q \Bigr) .
  \end{aligned}
  \label{eq:17.111}
\end{equation}
Passing to a continuous evolution, we find that the function $f_\gamma(x, Q)$
is determined by the integral-differential equation
\begin{equation}
  \frac{df_\gamma(x, Q)}{d\log Q} = \frac{\alpha}{2\pi}\ \int_x^1 \frac{dz}{z}\
    \frac{1+(1-z)^2}{z}\ f_e\Bigl( \frac{x}{z}, Q \Bigr) .
  \label{eq:17.112}
\end{equation}
Similarly, the distribution of component electrons in the physical electron
will evolve with $Q$, reflecting the appearance of electrons at lower values
of $x$ due to photon radiation, and the disappearance of these electrons at
higher $x$.  The term in brackets in Eq.~\eqref{eq:17.107} gives a correct
accounting of both effects.

\subsection{The Evolution Equations}

Collecting all of the splitting processes, we can write a complete set of
coupled evolution equations for the electron, positron, and photon
distribution functions in QED:
\begin{equation}
  \frac{df_f(x,Q)}{d\log Q} = \frac{\alpha}{2\pi}\ \int_x^1 \frac{dz}{z}\
    \sum_{f'} \bigl[ P_{f\to f'}(z)\,f_{f'}\Bigl( \frac{x}{z}, Q \Bigr)
    \bigr] ,
  \label{eq:17.120}
\end{equation}
where the splitting functions $P_{f\to f'}(z)$ are given by
\begin{equation}
  \begin{aligned}
    P_{e\to e\gamma}(z) &= \frac{1+z^2}{1-z},\\
    P_{e\to \gamma e}(z) &= 1 + (1-z)^2,\\
    P_{\gamma\to e\bar{e}}(z) &= z^2 + (1-z)^2,\\
    P_{e\to e}(z) &= \frac{3}{2}\,\delta(1-z) .
  \end{aligned}
  \label{eq:17.121}
\end{equation}
To obtain the distribution functions for an electron relevant to a given
momentum transfer $Q$, we should integrate these equations with the initial
conditions
\begin{equation}
  f_e(x,Q) = \delta(1-x),\qquad f_{\bar{e}}(x,Q) = 0,\qquad
  f_\gamma(x,Q) = 0 ,
  \label{eq:17.122}
\end{equation}
at $Q = m$.  With different initial conditions, the same equations give the
distribution functions for a physical positron or photon.  The solutions to
these equations are used as in Eq.~(17.115) to compute cross
sections involving processes induced by electrons, positrons, or photons that
involve large momentum transfer.

The evolution equations \eqref{eq:17.120} are constructed in such a way as to
conserve electron number and longitudinal momentum.  Thus, the basic sum
rules \eqref{eq:17.36} and \eqref{eq:17.39} satisfied by the parton
distributions of hadrons also apply to the QED distribution functions.
Specifically, the distribution functions of the electron contain one net
electron constituent,
\begin{equation}
  \int_0^1 dx\ \bigl[ f_e(x,Q) - f_{\bar{e}}(x,Q) \bigr] = 1 ,
  \label{eq:17.123}
\end{equation}
and account for the total momentum of the physical electron,
\begin{equation}
  \int_0^1 dx\ x\ \bigl[ f_e(x,Q) + f_{\bar{e}}(x,Q) + f_\gamma(x,Q) \bigr]
  = 1 .
  \label{eq:17.124}
\end{equation}
It is an instructive exercise to verify explicitly, using Eq.~\eqref{eq:17.120},
that the values of these integrals do not depend on $Q$.

\subsection{The Altarelli-Parisi Equations}

If we encounter mass singularities in QED associated with collinear photon
emission, we must also encounter mass singularities in QCD associated with
collinear gluon and quark emission.  If we compute the corrections of order
$\alpha_s$ to the leading-order parton cross sections discussed in
Sections~17.3 and~17.4, using massless quarks and gluons, we will find that
these correction terms diverge when we integrate over the collinear
configurations.  Thus the parton-model expressions, at least in their
simplest form, break down already at the next-to-leading order in
$\alpha_s$.

However, assuming that the singularities of QCD are no worse than those of
QED, the considerations of the previous section tell us how to treat these
singular terms.  In QED, we found it natural to include the large corrections
associated with the mass singularities in the parton distributions rather
than in the hard-scattering cross sections.  Viewed in this way, the singular
terms supply the kernel of an evolution equation for the parton distributions
as a function of the logarithm of the momentum scale.  Hard scattering with a
momentum transfer $Q$ probes the electron at a distance of order $Q^{-1}$.
When the electron wavefunction is resolved to very small scales, it appears as
a constituent electron, carrying only a fraction of the total longitudinal
momentum, plus a number of constituent photons and electron-positron pairs.
Any one of these constituents that carries a substantial fraction of the
total electron momentum can initiate a hard-scattering process.

Precisely the same logic applies to the calculation of QCD cross sections.
The contributions from the region of collinear gluon or quark emission should
be associated with the parton distribution functions rather than with the
hard-scattering cross sections.  If we make this association, we find that
the parton distributions are no longer independent of the momentum $Q$ that
characterizes the hard-scattering process; rather, they now evolve
logarithmically with $Q$.  For example, the basic equation \eqref{eq:17.30}
for deep-inelastic scattering will become
\begin{equation}
  \frac{d^2\sigma}{dx\,dQ^2}\bigl( e^{-}p \to e^{-}X \bigr)
  = \frac{2\pi\alpha^2}{Q^4}\ \Bigl[ 1 + (1-y)^2 \Bigr]\ \sum_f Q_f^2\,
    f_f(x, Q) ,
  \label{eq:17.125}
\end{equation}
and so Bjorken scaling will be violated.  Since this violation takes place
only on a logarithmic scale in $Q^2$, it will be a subtle effect, and
approximate Bjorken scaling will still be a prediction of QCD.  But the
violation of Bjorken scaling is inevitable, since QCD is a quantum field
theory with degrees of freedom at all momentum scales.  As we probe the
proton wavefunction at increasingly short distances, we excite the
high-momentum degrees of freedom and resolve the wavefunction into an
increasing number of quarks, antiquarks, and gluons.

The evolution of the QED parton distributions, governed by Eq.~\eqref{eq:17.120},
is characterized by the parameter $\alpha/\pi$, so the parton distributions
change by $\sim 1\%$ as $Q$ is changed by a factor of 10.  In QCD, the
corresponding factor governing the rate of evolution should be
$\alpha_s(Q)/\pi$.  Thus, when $Q$ is very small, the evolution is rapid and
contributions of higher order in perturbation theory are important.
Ultimately, the initial conditions for the evolution are determined by the
form of the proton wavefunction at large distance scales, which cannot be
calculated using Feynman diagrams.  On the other hand, when $Q$ is large,
well above 1 GeV in practice, the evolution becomes slow and is dominated by
the leading order in perturbation theory.  In that case, QCD perturbation
theory makes precise predictions for the form of the evolution of the parton
distributions, and these predictions can be tested against experiment.

To derive the evolution equations of parton distributions in QCD we can use
the same techniques and logic that we used above for QED.  There is a
subtlety, that the reduction of the gluon propagator to transverse
polarization states in the limit $q^2\to 0$, Eq.~\eqref{eq:17.81}, cannot be
proved so simply as in QED.  However, the result is correct also in the
non-Abelian case.%
\footnote{See, for example, J.~Collins and D.~Soper, in A.~Mueller,
\emph{Quantum Chromodynamics} (World Scientific, Singapore, 1991).}  Once
this technical point is resolved, the kinematics of collinear emission is
exactly the same as in QED.  Thus we find evolution equations of the same
form as in QED, modified only by the replacement of $\alpha$ by $\alpha_s$,
the insertion of appropriate color factors, and the accounting of the effects
of the three-gluon vertex.

\begin{figure}[t]
  \centering
  \includegraphics[width=0.6\textwidth]{fig17.20.png}
  \caption{The three vertices that contribute to parton evolution in QCD.}
  \label{fig:17.20}
\end{figure}

Collinear emission processes in QCD involve the three vertices shown in
Fig.~\ref{fig:17.20}.  Of these, the first two have the same Lorentz structure
as those shown in Figs.~17.16 and 17.19.  The only
difference, aside from the strength of the coupling constant, comes in the
color indices.  We will treat color just as we treated spin in the preceding
analysis: We average over initial colors, and sum over final colors.  Then
the first vertex of Fig.~\ref{fig:17.20}, representing the splitting of a
quark into a quark and a gluon, receives the color factor
\begin{equation}
  \frac{1}{2}\mathrm{tr}[t^at^a] = C_2(r) = \frac{4}{3} .
  \label{eq:17.126}
\end{equation}
The second vertex, representing the splitting of a gluon into a
quark-antiquark pair, receives the factor
\begin{equation}
  \frac{1}{2}\mathrm{tr}[t^at^a] = \frac{1}{2} .
  \label{eq:17.127}
\end{equation}
The third vertex in Fig.~\ref{fig:17.20} represents the splitting of a gluon
to two gluons, an effect that is new to the non-Abelian case.  It is
straightforward to compute the contribution of this vertex to the evolution
equations by taking the matrix elements of the vertex between transverse
gluon states of definite helicity.  This calculation is the subject of
Problem~17.4.

By accounting for all of these effects, we can modify the QED evolution
equations \eqref{eq:17.120} into the correct set of evolution equations for
parton distributions in QCD.  These are known as the \emph{Altarelli-Parisi
equations}.  They describe the coupled evolution of parton distributions
$f_q(x,Q)$, $f_{\bar{q}}(x,Q)$ for each flavor of quark and antiquark that
can be treated as massless at the scale $Q$, together with the parton
distribution of gluons, $f_g(x,Q)$.  Explicitly,
\begin{equation}
  \begin{aligned}
    \frac{df_g(x,Q)}{d\log Q} &= \frac{\alpha_s(Q^2)}{\pi}\ \int_x^1
      \frac{dz}{z}\ \Bigl[ P_{g\to g}(z)\,f_g\Bigl( \frac{x}{z}, Q \Bigr)
      + P_{g\to q}(z)\sum_f f_f\Bigl( \frac{x}{z}, Q \Bigr) \Bigr],\\
    \frac{df_f(x,Q)}{d\log Q} &= \frac{\alpha_s(Q^2)}{\pi}\ \int_x^1
      \frac{dz}{z}\ \Bigl[ P_{q\to q}(z)\,f_f\Bigl( \frac{x}{z}, Q \Bigr)
      + P_{q\to g}(z)\,f_g\Bigl( \frac{x}{z}, Q \Bigr) \Bigr] .
  \end{aligned}
  \label{eq:17.128}
\end{equation}
The first three splitting functions can be taken from Eqs.~\eqref{eq:17.121},
multiplied by the color factors computed in Eqs.~\eqref{eq:17.126} and
\eqref{eq:17.127}:
\begin{equation}
  \begin{aligned}
    P_{q\to qg}(z) &= \frac{4}{3}\ \frac{1+z^2}{1-z},\\
    P_{q\to gq}(z) &= \frac{4}{3}\ \frac{1+(1-z)^2}{z},\\
    P_{g\to q\bar{q}}(z) &= \frac{1}{2}\ \bigl[ z^2 + (1-z)^2 \bigr] .
  \end{aligned}
  \label{eq:17.129}
\end{equation}
The fourth splitting function requires also the computation of Problem~17.4;
the result is
\begin{equation}
  P_{g\to g}(z) = 6\,\delta(1-z) + 6\ \frac{z}{(1-z)_+}
    + 6z(1-z) - \frac{2}{3}n_f\,\delta(1-z) .
  \label{eq:17.130}
\end{equation}
The final term in this expression, which is proportional to $n_f$, the number
of light quark flavors, is the subtraction term associated with gluon
splitting into $q\bar{q}$ pairs.  The Altarelli-Parisi equations describe the
evolution of parton distributions for any hadron, or any hadronic
constituent, up to corrections of order $\alpha_s$ that are not enhanced by
large logarithms.

Our derivation of the Altarelli-Parisi equations respects the conservation
laws of QCD for quark numbers and longitudinal momentum.  Thus, the equations
must respect the parton-model sum rules \eqref{eq:17.36} and
\eqref{eq:17.39}.  As in the QED case, it is instructive to verify
explicitly that these integrals are independent of $Q$.

\begin{figure}[t]
  \centering
  \includegraphics[width=0.6\textwidth]{fig17.21.png}
  \caption{The $u$ quark parton distribution function $xf_u(x,Q)$ at
    $Q = 2$, 20, and 200 GeV, showing the effects of parton evolution
    according to the Altarelli-Parisi equations.  These curves are taken from
    the CTEQ fit to deep inelastic scattering data described in
    Fig.~\ref{fig:17.6}.}
  \label{fig:17.21}
\end{figure}

In QED, we could use the evolution equations to explicitly compute the
structure function of the electron.  In QCD, this is no longer possible,
because the initial conditions required to integrate the equations are
determined by the strong-coupling region of QCD and so are not known a
priori.  However, one can determine the initial conditions of the proton
structure experimentally, by measuring the cross section for deep inelastic
scattering at a given value of $Q^2$.  One can then predict the structure
functions, and thus the deep inelastic cross sections, at higher values of
$Q^2$.  There is one subtlety in this analysis: The gluon distribution is not
directly measured in deep inelastic scattering, but it does enter the
evolution equation for the quark distributions.  Thus, some of the information
on the $Q^2$ dependence of deep inelastic scattering simply goes into
determining the gluon distribution.  However, the gluon distribution is
absolutely normalized by the momentum sum rule \eqref{eq:17.39}, so the
evolution equations have predictive power even if this distribution must be
fit from the data.

The Altarelli-Parisi equations predict a characteristic form for the
evolution of parton distribution functions, shown in Fig.~\ref{fig:17.21}.
Partons at high $x$ tend to radiate and drop down to lower values of $x$.
Meanwhile, new partons are formed at low $x$ as products of this radiation.
Thus, the parton distributions decrease at large $x$ and increase much more
rapidly at small $x$ as $Q^2$ increases.  We can picture the proton as having
more and more constituents, but each carrying a smaller fraction of the total
momentum, as it is probed at smaller and smaller distances.

\section{Measurements of $\alpha_s$}

We have now developed the basic tools of perturbative QCD, and have applied
them to a number of hard-scattering processes.  In this section we will
review the experimental determinations of $\alpha_s$, the strong coupling
constant, which test the quantitative predictions of QCD.  The values of
$\alpha_s$ extracted from various measurements are collected in
Table~\ref{tab:17.1}.  The table shows the values of $\alpha_s$ extracted
from each of these measurements, expressed in terms of the value in the
reference conventions, $\alpha_s^{\overline{MS}}(m_Z^2)$.

\begin{table}[t]
  \centering
  \caption{Values of $\alpha_s(m_Z)$ Obtained from QCD Experiments}
  \begin{tabular}{lll}
    \toprule
    Process & $\alpha_s(m_Z)$ & $Q$ (GeV)\\
    \midrule
    Deep inelastic scattering & 0.118 (6) & 1.7\\
    $R$ in $\tau$ lepton decay & 0.123 (4) & 1.8\\
    $J/\psi$, $\Upsilon$ spectroscopy & 0.110 (6) & 2.3\\
    Transverse momentum of $W$ production & 0.121 (24) & 4.\\
    Deep inelastic scattering (evolution) & 0.112 (4) & 5.\\
    Event shapes in $e^{+}e^{-}$ annihilation & 0.121 (6) & 5.8, 9.1\\
    Rate for $J/\psi$, $\Upsilon$ decay & 0.108 (10) & 9.5\\
    $R$ in $e^{+}e^{-}$ annihilation (20--65 GeV) & 0.124 (21) & 35.\\
    $R$ in $Z^0$ decay & 0.124 (7) & 91.2\\
    \bottomrule
  \end{tabular}
  \label{tab:17.1}
\end{table}

The values of $\alpha_s(m_Z)$ displayed in this table are obtained by fitting
experimental results to the theoretical expressions given by perturbative QCD
using minimal subtraction.  The values of $\alpha_s$ have been evolved to
$Q = m_Z$ using the renormalization group equation.  $R$ refers to the ratio
of cross sections or partial widths to hadrons versus leptons.  The numbers
in parentheses are the standard errors in the last displayed digits.  The
column labelled `$Q$' gives an idea of the value of $Q$ at which the
measurement was made.  (Typically, these measurements average over a range
of $Q$, and that averaging is taken into account in the quoted values of
$\alpha_s$.)  This table is based on the results compiled by I.~Hinchliffe in
his article for the Particle Data Group, \emph{Phys.\ Rev.\ D} \textbf{50},
1297 (1994).  This article contains a full set of references and a discussion
of the sources of uncertainty in these determinations.

We see that several of the experiments determine $\alpha_s$ to an accuracy of
5\%, and that the various determinations are consistent with one another at
this level.  In Fig.~\ref{fig:17.23}, we have plotted the original values of
$\alpha_s$ represented in Table~\ref{tab:17.1}, before conversion to a common
scale, versus the momentum scale $Q$ at which each was obtained.  This
comparison gives a striking direct verification of the running of
$\alpha_s$.

\begin{figure}[t]
  \centering
  \includegraphics[width=0.6\textwidth]{fig17.23.png}
  \caption{Measurements of $\alpha_s$, plotted against the momentum scale $Q$
    at which the measurement was made.  This figure was constructed by
    evolving the values of $\alpha_s(m_Z)$ listed in Table~\ref{tab:17.1} back
    to the values of $Q$ indicated in the table.  The value for $e^{+}e^{-}$
    event shapes has been split into two points corresponding to experiments
    at the TRISTAN and LEP accelerators.  These values are compared to the
    theoretical expectation from the renormalization group evolution with the
    initial condition $\alpha_s(m_Z) = 0.117$.}
  \label{fig:17.23}
\end{figure}

At the beginning of this chapter, we wrote down a candidate for the
fundamental theory of strong interactions using only a few simple principles:
the existence of quarks and the identification of their quantum numbers, and
the idea that the theory of the quark interactions should be an
asymptotically free gauge theory.  It is remarkable that these simple
considerations have led us to a description of strong interactions that is
quantitatively correct for a broad range of phenomena in the hard-scattering
regime where asymptotic freedom can be used as a tool for calculation.

\problems

\textbf{17.1} \emph{Two-loop renormalization group relations.}
\begin{enumerate}[label=(\alph*)]
  \item In higher orders of perturbation theory, the expression for the QCD
    $\beta$ function will be a series
    \begin{equation}
      \beta(g) = -b_0\frac{g^3}{16\pi^2} - b_1\frac{g^5}{(16\pi^2)^2}
        + \cdots .
    \end{equation}
    Integrate the renormalization group equation and show that the running
    coupling constant is now given by
    \begin{equation}
      \alpha_s(Q) \approx \frac{4\pi}{b_0}\ \Bigl[ \frac{1}{\log(Q^2/\Lambda^2)}
        - \frac{b_1}{b_0}\ \frac{\log(\log(Q^2/\Lambda^2))}
          {(\log(Q^2/\Lambda^2))^2} + \cdots \Bigr],
    \end{equation}
    where the omitted terms decrease as $(\log(Q^2/\Lambda^2))^{-2}$.
  \item Combine this formula with the perturbation series for the
    $e^{+}e^{-}$ annihilation cross section:
    \begin{equation}
      \sigma(e^{+}e^{-}\to\text{hadrons}) = \sigma_0\cdot 3\cdot
        \sum_f Q_f^2\ \Bigl[ 1 + a_1\alpha_s + a_2\alpha_s^2
          + \mathcal{O}(\alpha_s^3) \Bigr] .
    \end{equation}
    The coefficient $a_2$ depends on the details of the renormalization
    conditions defining $\alpha_s$.  Show that the leading two terms in the
    asymptotic behavior of $\sigma(s)$ for large $s$ depend only on $b_0$ and
    $b_1$ and are independent of $a_2$ and $b_2$.  Thus the first two
    coefficients of the QCD $\beta$ function are independent of the
    renormalization prescription.
\end{enumerate}

\textbf{17.2} \emph{A direct test of the spin of the gluon.}  In this
problem, we compare the predictions of QCD with those of a model in which the
interaction of quarks is mediated by a scalar boson.  Let the coupling of the
scalar gluon to quarks be given by
\begin{equation}
  \delta\mathcal{L} = g_S\bar{q}q\,S ,
\end{equation}
and define $\alpha_g = g^2/4\pi$.
\begin{enumerate}[label=(\alph*)]
  \item Using the technique described in parts (b) and (c) of the Final
    Project of Part~I, compute the cross section for $e^{+}e^{-}\to q\bar{q}S$
    to the leading order of perturbation theory.  This cross section depends
    on the energies of the $q$, $\bar{q}$, and $S$, which we represent as
    fractions $x_1$, $x_2$, $x_3$ of the electron beam energy, as in
    Eq.~\eqref{eq:17.18}.  Show that
    \begin{equation}
      \frac{d^2\sigma}{dx_1\,dx_2}(e^{+}e^{-}\to q\bar{q}S)
      = \frac{4\pi\alpha^2}{3s}\cdot 3Q_f^2\cdot
        \frac{\alpha_g}{\pi}\ \frac{x_2}{(1-x_1)(1-x_3)^2}\ \cdots .
    \end{equation}
  \item In practice, it is very difficult to tell quarks from gluons
    experimentally, since both particles appear as jets of hadrons.
    Therefore, let $x_a$ be the largest of $x_1$, $x_2$, $x_3$, let $x_b$ be
    the second largest, and let $x_c$ be the smallest.  Sum over the various
    possibilities to derive an expression for $d^2\sigma/dx_adx_b$, both in
    QCD, using Eq.~\eqref{eq:17.18}, and in the scalar gluon model.  Show
    that these models can be distinguished by their distributions in the
    $(x_a, x_b)$ plane.
\end{enumerate}

\textbf{17.3} \emph{Quark-gluon and gluon-gluon scattering.}
\begin{enumerate}[label=(\alph*)]
  \item Compute the differential cross section
    \begin{equation}
      \frac{d\hat{\sigma}}{d\hat{t}}\bigl( q\bar{q} \to gg \bigr)
    \end{equation}
    for quark-antiquark annihilation in QCD to the leading order in
    $\alpha_s$.  This is most easily done by computing the amplitudes between
    states of definite quark and gluon helicity.  Ignore all masses.  Use
    explicit polarization vectors and spinors, for example,
    \begin{equation}
      \epsilon^\mu = \frac{1}{\sqrt{2}}(0, 1, i, 0)
    \end{equation}
    for a right-handed gluon moving in the $+3$ direction.  You need only
    consider transversely polarized gluons.  By helicity conservation, only
    the initial states $q_L\bar{q}_R$ and $q_R\bar{q}_L$ can contribute; by
    parity, these two states give identical cross sections.  Thus it is
    necessary only to compute the amplitudes for the three processes
    \begin{equation}
      q_L\bar{q}_R \to g_Rg_R;\qquad
      q_L\bar{q}_R \to g_Rg_L;\qquad
      q_L\bar{q}_R \to g_Lg_L .
    \end{equation}
    In fact, by $CP$ invariance, the first and third processes have equal
    cross sections.  After computing the amplitudes, square them and combine
    them properly with color factors to construct the various helicity cross
    sections.  Finally, combine these to form the total cross section
    averaged over initial spins and colors.
  \item Compute the differential cross section
    \begin{equation}
      \frac{d\hat{\sigma}}{d\hat{t}}\bigl( gg \to gg \bigr)
    \end{equation}
    for gluon-gluon scattering.  There are 16 possible combinations of
    helicities, but many of them are related to each other by parity and
    crossing symmetry.  All 16 can be built up from the three amplitudes for
    \begin{equation}
      g_Rg_R \to g_Rg_R;\qquad
      g_Rg_R \to g_Rg_L;\qquad
      g_Rg_R \to g_Lg_L .
    \end{equation}
    Show that the last two of these amplitudes vanish.  The first can be
    dramatically simplified using the Jacobi identity.  When the smoke
    clears, only three of the 16 polarized gluon scattering cross sections
    are nonzero.  Combine these to compute the spin- and color-averaged
    differential cross section.
\end{enumerate}

\textbf{17.4} \emph{The gluon splitting function.}  Compute the gluon
splitting function \eqref{eq:17.130} for the Altarelli-Parisi equations.  To
carry out this computation, first compute the matrix elements of the
three-gluon vertex shown in Fig.~\ref{fig:17.20} between gluon states of
definite helicity.  Combine these to derive the splitting function in the
region $x < 1$.  Then fix the singularity of the splitting function at
$x = 1$ to give this function the correct overall normalization.

\textbf{17.5} \emph{Photoproduction of heavy quarks.}  Consider the process of
heavy quark pair photoproduction, $\gamma + p \to Q\bar{Q} + X$, for a heavy
quark of mass $M$ and electric charge $Q$.  If $M$ is large enough, any
diagram contributing to this process must involve a large momentum transfer;
thus a perturbative QCD analysis should apply.  This idea applies in practice
already for the production of $c$ quark pairs.  Work out the cross section to
the leading order in QCD.  Choose the parton subprocess that gives the
leading contribution to this reaction, and write the parton-model expression
for the cross section.  You will need to compute the relevant subprocess
cross section, but this can be taken directly from one of the QED
calculations in Chapter~5.  Then use this result to write an expression for
the cross section for $\gamma$-proton scattering.

\textbf{17.6} \emph{Behavior of parton distribution functions at small $x$.}
It is possible to solve the Altarelli-Parisi equations analytically for very
small $x$, using some physically motivated approximations.  In this limit,
the gluon distribution dominates, and the evolution is driven by the
singularity of $P_{g\to g}(z)$ at small $z$.  Show that the gluon distribution
grows as a power of $x$ as $x\to 0$, and derive the power in terms of
$\alpha_s$.
